\documentclass[12pt,a4paper]{article}
\pdfvariable trailerid {[
  <D3A301092026D3A301092026D3A30109>
  <D3A301092026D3A301092026D3A30109>
]}
\pdfextension catalog {/Lang (en-US)}

\usepackage[english]{babel}
\usepackage{amsmath}
\usepackage{amssymb}
\usepackage{amsthm}
\usepackage{mathtools}
\usepackage{mathrsfs}
\usepackage{microtype}
\usepackage[margin=2.6cm]{geometry}
\usepackage{booktabs}
\usepackage{array}
\usepackage{longtable}
\usepackage{xurl}
\usepackage[hidelinks]{hyperref}
\Urlmuskip=0mu plus 2mu\relax
\urlstyle{tt}
\hypersetup{
  pdftitle={Depth-Three Smith Profiles and Newton-Minor Geometry for Affine Hjelmslev-Radon Incidence},
  pdfauthor={Oleksiy Babanskyy},
  pdfsubject={Integral Smith theory of depth-three affine Hjelmslev-Radon incidence},
  pdfkeywords={Smith normal form, Hjelmslev geometry, Radon transform, q-Pascal matrix, Newton minors}
}
\usepackage[nameinlink,noabbrev]{cleveref}
\usepackage{csquotes}
\usepackage[backend=biber,style=numeric-comp,sorting=nyt,maxbibnames=12,giveninits=true]{biblatex}
\addbibresource{references.bib}
\setcounter{biburlnumpenalty}{9000}
\setcounter{biburlucpenalty}{9000}
\setcounter{biburllcpenalty}{9000}
\AtBeginBibliography{\small\raggedright
  \setlength{\emergencystretch}{3em}%
  \setlength{\bibitemsep}{0pt}%
  \hbadness=10000}

\newtheorem{theorem}{Theorem}[section]
\newtheorem{lemma}[theorem]{Lemma}
\newtheorem{proposition}[theorem]{Proposition}
\newtheorem{corollary}[theorem]{Corollary}
\theoremstyle{definition}
\newtheorem{definition}[theorem]{Definition}
\newtheorem{remark}[theorem]{Remark}

\newcommand{\Z}{\mathbb Z}
\newcommand{\Q}{\mathbb Q}
\newcommand{\F}{\mathbb F}
\newcommand{\R}{\mathsf R}
\newcommand{\A}{\mathsf A}
\newcommand{\OO}{\mathcal O}
\newcommand{\LL}{\mathsf L}
\newcommand{\QQ}{\mathsf Q}
\newcommand{\Rel}{\operatorname{Rel}}
\newcommand{\coker}{\operatorname{coker}}
\newcommand{\im}{\operatorname{im}}
\newcommand{\rank}{\operatorname{rank}}
\newcommand{\length}{\operatorname{length}}
\newcommand{\ord}{\operatorname{ord}}
\newcommand{\Ann}{\operatorname{Ann}}
\newcommand{\Fib}{\operatorname{Fib}}
\newcommand{\ngen}{\operatorname{ngen}}
\newcommand{\Hilb}{\operatorname{Hilb}}
\newcommand{\vp}{v_p}
\newcommand{\eps}{\varepsilon}
\newcommand{\id}{\mathrm{id}}
\newcommand{\transpose}{\mathsf T}
\newcommand{\pos}[1]{\left[#1\right]_{+}}
% Generated deterministically from companion/formulas.json; do not edit by hand.
% FCIG-FORMULAS-SHA256: E307305E6DC754C0DB06DA4E0477332050258B3C38F07F4CD151475BE9B652B7
\newcommand{\FCIGDThreeEpsilonClosedFormula}{%
\begin{cases}
 18,&p=3,\\[1mm]
 \frac{(p-1)^2(4p^3+5p^2+2p+4)}{36},&p>3,\ p\equiv1\pmod6,\\[3mm]
 \frac{p(p+1)(8p^3-18p^2+11p-5)}{72},&p\equiv5\pmod6.
\end{cases}%
}
\newcommand{\FCIGDThreeSmithProfileTwo}{(30,6,19,6,2,1)}
\newcommand{\FCIGDThreeSmithProfileThree}{(240,168,195,108,15,3)}
\newcommand{\FCIGDThreeModularRankOne}{\frac{p(p+1)(3p^4+4p^3+3p-1)}{18}}
\newcommand{\FCIGDThreeModularRankTwo}{\frac{p^2(p+1)^2(3p^2+p+1)}{18}}
\newcommand{\FCIGDThreeSmithProfileModOne}{%
\boxed{\begin{aligned}
 m_0 &=\frac{p(p+1)(3p^4+4p^3+3p-1)}{18},\\
 m_1 &=\frac{p(p-1)^2(p+1)(6p^2+2p+1)}{18},\\
 m_2 &=\frac{p(p-1)(2p+1)(6p^3-p^2+9p-2)}{36},\\
 m_3 &=\frac{p(p-1)^2(p+1)(3p^2+4p-1)}{18},\\
 m_4 &=\frac{p(p-1)^2(p+2)}{4},\\
 m_5 &=\frac{p(p-1)}{2}.
\end{aligned}}%
}
\newcommand{\FCIGDThreeSmithProfileModFive}{%
\boxed{\begin{aligned}
 m_0 &=\frac{p^2(p+1)^2(3p^2+p+1)}{18},\\
 m_1 &=\frac{p(p+1)(2p-1)(6p^3-7p^2-5p+2)}{36},\\
 m_2 &=\frac{p(12p^5-8p^4+19p^3-12p^2-11p+4)}{36},\\
 m_3 &=\frac{p(p-1)(p+1)(6p^3+2p^2-11p+2)}{36},\\
 m_4 &=\frac{p(p-1)^2(p+2)}{4},\\
 m_5 &=\frac{p(p-1)}{2}.
\end{aligned}}%
}


\crefname{theorem}{theorem}{theorems}
\crefname{lemma}{lemma}{lemmas}
\crefname{proposition}{proposition}{propositions}
\crefname{corollary}{corollary}{corollaries}
\crefname{definition}{definition}{definitions}

\title{Depth-Three Smith Profiles and Newton-Minor Geometry for Affine
  Hjelmslev--Radon Incidence}
\author{Oleksiy Babanskyy}
\date{}
\pagestyle{plain}

\begin{document}
\raggedbottom
\maketitle

\begin{abstract}
Let $B_3$ be the integer line-by-point incidence matrix of the affine plane
over $\Z/p^3\Z$, with primitive normal directions taken modulo units.  We
determine every nonzero Smith factor for every prime.  The exponent is
$p^5$, and the six multiplicities $m_0,\ldots,m_5$ are given by explicit
polynomials on the two residue classes $p\equiv1,5\pmod6$, together with
separate exact rows for $p=2,3$.  The new ingredient is a reduction of the
first unresolved integral obstruction to a truncated $q$-Pascal common
image.  An exact $q$-Lucas block product identifies its supported columns;
a scalar recurrence proves module closure and boundary generation; and a
$q$-Newton row basis makes all determinantal ideals triangular with Smith
thresholds
\[
 \tau_p(K)=\mu_p\bigl(p^2(p-1)-K\bigr)+p\beta_p(K).
\]
This also yields the complete ambient filtered profile, Hilbert series and
minimal number of generators of the common image.  A typed cross-chart
exchange closes the remaining transfer defect.  Finite computations are
used only for exceptional certificates and regression, not to interpolate
the all-prime formulas.
\end{abstract}

\smallskip
\noindent\textbf{Keywords:} Smith normal form; affine Hjelmslev geometry;
Radon transform; incidence matrix; residue ring; $q$-Pascal matrix; Newton
minors.

\smallskip
\noindent\textbf{MSC 2020:} 05B20 (primary); 15B36, 51E15, 13H10
(secondary).

\begin{center}
\small
\href{https://github.com/aconsciousfractal/FCIG-Depth-Three-Smith-Profiles-and-Newton-Minor-Geometry-for-Affine-Hjelmslev-Radon-Incidence}{Companion repository and reproducibility materials}
\end{center}

\smallskip
\noindent\textcopyright\ 2026 Oleksiy Babanskyy.  The manuscript source
and compiled paper retain the author's copyright.  The accompanying original
software and repository documentation are released under the MIT License;
see \texttt{LICENSE\_SCOPE.md}.  No novelty, firstness, priority, or
exhaustive-literature claim is made.

\section{Contribution and logical spine}
\label{sec:contribution}

The incidence object in this paper is affine and integral.  It is neither a
projective incidence matrix nor a matrix over a field having the same
cardinality as $\Z/p^3\Z$.  The main result, \cref{thm:main-profile}, gives
the complete Smith profile at depth three for every prime.  Its proof is
not a finite-prime fit.  It combines four uniform layers:

\begin{enumerate}
\item a maximal-order and point-fibre transfer analysis, which fixes the
  total determinant order and high Smith tail;
\item a complete-ideal calculation of the modular rank;
\item an odd-prime tripotent-sector reduction to a truncated $q$-Pascal
  common image, followed by an exact Newton-minors calculation of that
  image;
\item a cross-chart divided-carry theorem, which removes the final transfer
  defect before the exact Smith bookkeeping is applied.
\end{enumerate}

The written proof and the finite calculations play distinct roles.
The exceptional primes $2$ and $3$ have exact finite certificates where a
semisimple or degree-bound argument does not apply.  The primes $5$, $7$,
and $11$ are retained as public exact regression rows evaluated from the
tracked formulas.  They do not enter the selection of either residue-class
branch, and none of these finite rows proves a uniform polynomial.

The diagonal cyclic valuation formulas used in the top-shell calculation
belong to Weber--K\"unzer \cite{WK04}; the cyclic CRT layer is also covered
by Mahatab--Sampath \cite{MS15}.  The complete-ideal correspondence and
Hoskin--Deligne length formula are imported through Lipman's work
\cite{LIP93,LW03}.  Sagan supplies $q$-Lucas \cite{SAG92}, while the
integral $q$-Stirling expansion is the Carlitz identity in the form recorded
by Cai--Ehrenborg--Readdy \cite{CER18}.  These sources do not state the
affine common-image theorem, the Newton-minors theorem, or the profile in
\cref{thm:main-profile}.

Sections~\ref{sec:object}--\ref{sec:modrank} fix the structural invariants.
Sections~\ref{sec:sector}--\ref{sec:filtered} compute the first integral
obstruction.  Section~\ref{sec:o2} closes the cross-chart transfer defect,
  and Sections~\ref{sec:assembly}--\ref{sec:controls} assemble and check the
answer.  Section~\ref{sec:repro} defines the public reproducibility and
repository boundary.

\section{Affine incidence lattice, maximal order,
\texorpdfstring{\protect\\}{ }and transfer}
\label{sec:object}

\begin{definition}[Affine Hjelmslev--Radon incidence]
\label[definition]{def:incidence}
Fix a prime $p$ and $n\geq1$, and put
\[
 \R_n=\Z/p^n\Z,\qquad G_n=\R_n^2.
\]
A vector $u=(a,b)$ is primitive if $a\R_n+b\R_n=\R_n$.  Primitive vectors
are identified under multiplication by $\R_n^\times$, and we use the chart
\[
 \mathbb P^1(\R_n)=
 \{(1,t):t\in\R_n\}\sqcup
 \{(ps,1):s\in\Z/p^{n-1}\Z\}.
\]
For $u\in\mathbb P^1(\R_n)$ and $c\in\R_n$, let
$\ell_{u,c}=\{x\in G_n:u\cdot x=c\}$.  The matrix $B_n$ has rows
$(u,c)$, columns $x\in G_n$, and entries
\[
 B_n[(u,c),x]=\mathbf 1_{\{u\cdot x=c\}}.
\]
Thus $B_n$ has $(p+1)p^{2n-1}$ rows, $p^{2n}$ columns, and each row has
$p^n$ ones.
\end{definition}

Let $\A_n=\Z[G_n]$.  For $H_u=\ker(u)$ put
$N_{H_u}=\sum_{h\in H_u}[h]$ and
\[
 \LL_n=\sum_u\A_nN_{H_u},\qquad \QQ_n=\A_n/\LL_n.
\]

\begin{lemma}[Group-ring dictionary]
\label[lemma]{lem:dictionary}
There is a canonical isomorphism
\[
 \QQ_n\cong\coker(B_n^{\transpose}).
\]
In particular, $\QQ_n$ records the $p^{2n}$ nonzero Smith factors of
$B_n$.
\end{lemma}

\begin{proof}
For fixed $u$, the $p^n$ affine fibres are disjoint translates of $H_u$;
their indicators form a $\Z$-basis of $\A_nN_{H_u}$.  Summing over the
direction chart identifies the image of $B_n^{\transpose}$ with $\LL_n$.
\end{proof}

The rational characters of exact order $p^k$ form Galois orbits indexed by
cyclic subgroups of order $p^k$.  These subgroups form a rooted projective
character tree; level $k$ has $(p+1)p^{k-1}$ nodes, each carrying
$\Z[\zeta_{p^k}]$.

\begin{theorem}[Maximal-order quotient]
\label[theorem]{thm:maximal-order}
Let $\OO_n$ be the maximal order in $\Q[G_n]$.  Then
\begin{align}
 \OO_n/\LL_n
 &\cong \Z/p^{2n}\Z\oplus
 \bigoplus_{k=1}^{n}
 (\Z/p^{2n-k}\Z)^{p^{2k}-p^{2k-2}},\tag{2.1}\label{eq:O-L}\\
 \vp|\QQ_n|
 &=\frac{np^{2n}}2-
   \frac{p(p^{2n}-1)}{2(p^2-1)}.\tag{2.2}\label{eq:Q-order}
\end{align}
In particular,
\[
 \vp|\QQ_3|=\frac{3p^6-p^5-p^3-p}{2}.
\]
\end{theorem}

\begin{proof}
For a leaf $C$ in the character tree let $\mathsf R_C$ be the image of the
cyclic group ring on the root-to-leaf components and put
$\mathsf M_n=\sum_C\mathsf R_C$.  Fourier evaluation of a line norm gives
the exact additive identity $\LL_n=p^n\mathsf M_n$.  Projection to the top
cyclotomic level is onto.  On each branch its kernel is generated by
$\Phi_{p^n}(t)$, whose value on every lower $p$-power root is $p$.
Top components are leaf-disjoint, hence
\[
 \mathsf M_n\cap\OO_{<n}=p\mathsf M_{n-1}.
\]
Induction gives exponent $n$ once at the root and exponent $n-k$ with
multiplicity $p^{2k}-p^{2k-2}$ at level $k$ for
$\mathsf M_n\subset\OO_n$.  Scaling by $p^n$ proves \eqref{eq:O-L}.
Subtracting the standard group-ring normalization index from the exponent
sum yields \eqref{eq:Q-order}.  The cyclic branch kernel is the prior-art
CRT layer \cite{MS15}; the projective-tree assembly is the step used here.
\end{proof}

Let $\pi_n:G_n\to G_{n-1}$ be reduction and define point-fibre transfer
\[
 \tau_n([a])=\sum_{\pi_n(x)=a}[x].
\]

\begin{theorem}[Exact transfer and multiscale pencil]
\label[theorem]{thm:transfer}
For $n\geq2$,
\[
 \tau_n(\A_{n-1})\cap\LL_n=\tau_n(\LL_{n-1}),
\]
so $\tau_n$ induces an injection $\QQ_{n-1}\hookrightarrow\QQ_n$.
Moreover, in $\QQ_n$,
\[
 p^n[x]=\tau_n\bigl(p^{n-2}[x\bmod p^{n-1}]\bigr),
 \qquad p^n\QQ_n=\tau_n(p^{n-2}\QQ_{n-1}).
\tag{2.3}
\]
Consequently $\exp(\QQ_n)=p^{2n-1}$ and, if $m_h^{(n)}$ denotes the
multiplicity of $\Z/p^h\Z$ in $\QQ_n$, then
\[
 m_{n+h}^{(n)}=m_{n-2+h}^{(n-1)}\qquad(1\leq h\leq n-1).
\tag{2.4}
\]
\end{theorem}

\begin{proof}
In Fourier coordinates a point fibre has value $p^2$ on lower characters
and zero on exact-order-$p^n$ characters.  Combining this with the
leafwise description in \cref{thm:maximal-order} gives
\[
 \tau_n(a)\in\LL_n\Longleftrightarrow a\in\LL_{n-1},
\]
and disjoint point-fibre supports give injectivity.

For the second assertion let $F_{r,n}(x)$ be the sum of points congruent to
$x$ modulo $p^r$, and let $S_x$ be the sum of the unique line of every
direction through $x$.  Coefficient counting by the valuation of $y-x$
gives in $\A_n$
\[
 S_x=F_{0,n}+\sum_{r=1}^{n-1}(p^r-p^{r-1})F_{r,n}(x)+p^n[x].
\]
Both $S_x$ and $F_{0,n}$ vanish in $\QQ_n$.  Subtracting the transferred
depth-$(n-1)$ identity leaves (2.3).  Injection and induction from the
depth-one affine profile \cite{CSX06} give the exact exponent; comparing
cyclic factors in (2.3) gives (2.4).
\end{proof}

At depth two write
\[
 \alpha=\frac{p^3(p-1)}2,\qquad
 \beta=\frac{p(p-1)^2(p+2)}4,\qquad
 \gamma=\frac{p(p-1)}2.
\tag{2.5}
\]
The depth-two profile has multiplicities $\alpha,\beta,\gamma$ at
exponents $1,2,3$ \cite{BAB26D2}.  Hence (2.1)--(2.4) imply for
depth three
\[
 m_4=\beta,\qquad m_5=\gamma,
\tag{2.6}
\]
and exclude exponents above five.

\section{The modular rank and the cyclic top-shell count}
\label{sec:modrank}

We next determine $m_0=\rank_{\F_p}B_3$.  Let
$\overline\A_3=\F_p[(\Z/p^3\Z)^2]$ and let $I_3$ be the span of all
affine-line indicators.  In augmentation coordinates
$\overline\A_3\cong k[x,y]/(x^{p^3},y^{p^3})$, put
$J_3=\Ann(I_3)$.

\begin{lemma}[Depth-three bundle annihilator]
\label[lemma]{lem:bundle-annihilator}
For each $e\in\mathbb P^1(\Z/p^2\Z)$ choose regular parameters
$(z_e,v_e)$ along the associated direction.  Then
\[
 J_3=\bigcap_e(z_e,v_e^{p^2})^p.
\tag{3.1}
\]
\end{lemma}

\begin{proof}
A primitive cyclic line subgroup $H=\langle h\rangle$ has norm
$N_H=(h-1)^{p^3-1}$ in characteristic $p$, hence
$\Ann(N_H)=(h-1)$.  Group the $p$ fine lifts above a depth-two direction
as $h_{e,a}=g_et_e^{p^2a}$.  With $z=g_e-1$ and $w=v^{p^2}$, their linear
factors are $(1+z)(1+w)^a-1$.  Modulo $(z,w)^p$, truncated logarithms send
  them to the $p$ distinct forms $Z+aW$.  A polynomial of degree below $p$
  divisible by all these forms is zero, so the intersection contains no
  nonzero class of degree below $p$.  Conversely, modulo each monic factor
  $(1+z)(1+w)^a-1$ one has $z\in(w)$; hence every monomial of total degree
  $p$ maps into $(w^p)=0$.  Thus $(z,w)^p$ lies in every factor ideal, and
  their intersection is exactly $(z,w)^p$.  Intersecting all bundles gives
  (3.1).
\end{proof}

\begin{lemma}[Formal lift and simultaneous base-point cluster]
\label[lemma]{lem:formal-cluster}
Put $N=p^3$, $q=p^2$, let $R=k[[x,y]]$ with maximal ideal
$\mathfrak m=(x,y)$, and choose formal lifts of the regular pairs in
\cref{lem:bundle-annihilator}.  Set
\[
\widetilde I_e=(z_e,v_e^q)^p,
\qquad
\widetilde J_3=\bigcap_{e\in\mathbb P^1(\Z/p^2\Z)}\widetilde I_e.
\]
Then $\mathfrak m^N\subseteq\widetilde J_3$, the ideal
$\widetilde J_3$ is complete, and reduction modulo $(x^N,y^N)$ identifies it
with $J_3$.  A common point-blowup resolution has
one root, $p+1$ trunks of length $p-1$, and $p$ branches above each trunk of
length
\[
 L=q-p=p(p-1).
\]
In the strict exceptional basis its componentwise-maximum divisor has
coefficients
\[
 p\quad\text{at the root},\qquad
 p(j+1)\quad(1\leq j\leq p-1)\quad\text{on a trunk},
\]
and
\[
 p(p+r)\quad(1\leq r\leq L)\quad\text{on a branch}.
\]
Its least antinef majorant represents $\widetilde J_3$.  In particular,
\[
 \rank_k B_3=\dim_k I_3
 =\dim_k(\overline\A_3/J_3)
 =\length_R(R/\widetilde J_3).
\]
\end{lemma}

\begin{proof}
In any regular pair $(z,v)$, a monomial $z^av^b$ outside
$(z,v^q)^p$ satisfies $a+\lfloor b/q\rfloor\leq p-1$, hence
$a+b\leq pq-1=N-1$.  Therefore $\mathfrak m^N$ lies in every
$\widetilde I_e$.  The same monomial description says
\[
 (z,v^q)^p=(z^av^b:a,b\geq0,\ qa+b\geq pq),
\]
whose exponent set is exactly the set of lattice points in its Newton
polyhedron.  Hence each $\widetilde I_e$ is integrally closed, and their
finite intersection is integrally closed as well: an element integral over
the intersection is integral over every $\widetilde I_e$, and hence belongs
to every $\widetilde I_e$.  Thus $\widetilde J_3$ is a complete
$\mathfrak m$-primary ideal.  Since
$(x^N,y^N)\subseteq\mathfrak m^N\subseteq\widetilde I_e$ for every $e$,
finite intersection commutes with the Artinian quotient and
\[
 \widetilde J_3/(x^N,y^N)=J_3,
 \qquad
 \length_R(R/\widetilde J_3)=\dim_k(\overline\A_3/J_3).
\]
Let $\bar k$ be an algebraic closure of $k$.  Because all ideals contain
$\mathfrak m^N$, the base change $R\to\bar k[[x,y]]$ may be computed in the
finite quotient $R/\mathfrak m^N$.  Faithful flatness makes it commute with
the finite intersection and gives
\[
 \dim_k(R/\widetilde J_3)
 =\dim_{\bar k}\bigl(\bar k[[x,y]]/
   \widetilde J_3\bar k[[x,y]]\bigr).
\]
We may therefore use the complete-ideal/antinef correspondence over
$\bar k$ and descend its numerical colength to $k$.

We now construct the common cluster.  In the finite-slope chart write a
direction as $a=a_0+pa_1$, with $a_0,a_1\in k$, and take
\[
 z_a=(1+x)^{-a}(1+y)-1,\qquad v_a=x.
\]
For two lifts of the same coarse direction, Frobenius gives
\[
 1+z_{a_0+pa_1}
 =(1+z_{a_0+pb_1})(1+x^p)^{\,b_1-a_1}.
\]
If $a_1\ne b_1$, reduction modulo $z_{a_0+pb_1}$ has first nonzero term
$(b_1-a_1)x^p$.  The two smooth branches therefore have intersection
multiplicity exactly $p$.  After $j$ blowups along their common smooth
branch, $0\leq j<p$, their difference has leading term a nonzero multiple of
$x^{p-j}$; after the $p$th it records the distinct residue $a_1$.  They
therefore share the root and $p-1$ successive free infinitely-near points,
and then separate at the $p$ distinct values of $a_1$.  Directions with
different $a_0$ have distinct linear tangents $y-a_0x$ and separate after
the first blowup.  Interchanging $x$ and $y$ treats the remaining chart
$(pa_1,1)$.  Hence there are exactly $p+1$ coarse trunks, with $p$ branches
over each trunk.  Every center lies on the strict transform of a smooth
branch and only on the newest exceptional component, so no center is
satellite.  The coarse tangents, the separation parameters $a_1$ and all
subsequent centers are $k$-rational; hence every residue-degree factor in the
Hoskin--Deligne sum is one.
For $p=2$ the same construction has three trunks of length one and two
branches of length two above each trunk; the coefficient $b_1-a_1$ remains
nonzero, so no contact or separation step degenerates.

For one direction, follow the smooth branch $z_e=0$.  At the $j$th center,
successive substitutions $z_{e,j-1}=v_ez_{e,j}$ give the weak transforms
\[
 (z_{e,j},v_e^{q-j})^p,\qquad 0\leq j<q.
\]
Thus $\widetilde I_e$ has point-basis multiplicity $p$ at the root and at
each of a free chain of $q-1$ further points.  Its first $p$ centers are the
root and the common trunk; after separation the remaining $q-p=L$ centers
form its branch.

Let $\mu:X\to\operatorname{Spec}R$ be the common resolution and let $D_e$
be the divisor of $\widetilde I_e$ on $X$.
Along its own free chain, point-basis multiplicity $p$ makes the coefficient
of the exceptional component created at depth $d$ in the strict exceptional
basis equal to $p(d+1)$.
An ideal whose branch has already separated has unit weak transform at a
center on another branch, so further pullback only propagates its coefficient
at the last common component; it cannot exceed the coefficient supplied by
the ideal belonging to that branch.
Taking $D^{\rm raw}=\max_eD_e$ componentwise therefore gives exactly the
displayed root, trunk and branch coefficients.  A rational function is a
local section of $\mathcal O(-D^{\rm raw})$ precisely when its exceptional
valuations dominate every $D_e$; hence
$\mathcal O(-D^{\rm raw})=\bigcap_e\mathcal O(-D_e)$ and
\[
 \mu_*\mathcal O(-D^{\rm raw})
 =\bigcap_e\mu_*\mathcal O(-D_e)
 =\bigcap_e\widetilde I_e=\widetilde J_3.
\]
Unloading replaces $D^{\rm raw}$ by its unique least antinef majorant
without changing this pushforward
\cite[Section~1]{LIP93}\cite[Corollary~1.1.2 and Lemma~1.2]{LW03}.
Finally, affine translation preserves the line span, so $I_3$ is an ideal.
The coefficient-of-identity pairing makes $\overline\A_3$ a symmetric
Frobenius algebra; its nondegeneracy therefore gives
$J_3=\Ann(I_3)=I_3^\perp$, and the final displayed equality follows.
\end{proof}

The preceding lemma supplies the raw divisor rather than assuming it.
Symmetry and unloading now give the following least antinef point basis.

\begin{theorem}[Depth-three modular rank]
\label[theorem]{thm:modrank}
Let $A$ be the root point-basis coefficient, $H$ the common point-basis
coefficient at every trunk point and $K_1,\ldots,K_L$ the successive
point-basis coefficients on each branch.  The least antinef majorant has
point basis
\[
\begin{array}{c|c|c|c}
\text{case}&A&H&(K_i)_{1\leq i\leq L}\\ \hline
p=3&15&3&1,\ldots,1\\
p=3h+1&(p+1)H+p&\dfrac{(p-1)(p+2)}3&h,\ldots,h\\[2mm]
p=3h+2&(p+1)H&\dfrac{p(p+1)}3&
 h+1\ (ph\text{ times}),\ h\ (L-ph\text{ times}).
\end{array}
\tag{3.2}
\]
The last line includes $p=2$.  Consequently
\[
 M_p:=\rank_{\F_p}B_3=
 \begin{cases}
 240,&p=3,\\[1mm]
 \FCIGDThreeModularRankOne,&p\equiv1\pmod3,\\[3mm]
 \FCIGDThreeModularRankTwo,&p\equiv2\pmod3.
 \end{cases}
\tag{3.3}
\]
\end{theorem}

\begin{proof}
The least antinef majorant is unique.  The decorated tree is invariant under
permutations of equal trunks and of the $p$ branches above each trunk, so the
least majorant has a symmetric point basis.  Write $h$ for its common trunk
coefficient at the first trunk point, and write $H_1,\ldots,H_{p-1}$ for
the successive trunk coefficients, so $H_1=h$.  On each branch its
point-basis coefficients satisfy
\[
 h=H_1\geq H_2\geq\cdots\geq H_{p-1}\geq pK_1,
 \qquad
 K_1\geq K_2\geq\cdots\geq K_L\geq0,
\]
while root proximity gives $A\geq(p+1)h$.  The corresponding strict
exceptional coefficients are
\[
 A+\sum_{q=1}^{j}H_q\quad(1\leq j\leq p-1),\qquad
 A+\sum_{q=1}^{p-1}H_q+\sum_{i=1}^{r}K_i
 \quad(1\leq r\leq L).
\]
The last one must dominate the raw endpoint coefficient $p^3$.  Since
$H_q\leq h$ and $K_i\leq\lfloor h/p\rfloor$, every symmetric antinef
majorant therefore satisfies
\[
 A\geq\max\left((p+1)h,
 p^3-(p-1)h-p(p-1)\left\lfloor\frac hp\right\rfloor\right).
\tag{3.4}
\]
Put $F_p(h)$ equal to the right-hand side.  Writing $h=pk+s$,
$0\leq s<p$, makes the second entry minus the first
\[
 p\bigl(p^2-(3p-1)k-2s\bigr).
\]
If $p=3m+1$, the unique minimum occurs at
$k=m$, $s=2m$, namely $h=(p-1)(p+2)/3$: the difference is $p$ there and
$-p$ at $h+1$, while the decreasing entry is larger by $p-1$ at $h-1$.
Thus $F_p(h)=(p+1)h+p$.  If $p=3m+2$, the unique minimum is at
$h=p(m+1)$; comparison with $h-1$ and $h+1$ gives $F_p(h)=(p+1)h$.
For $p=3$, direct evaluation gives the unique minimum $F_3(3)=15$.

It remains to show that the lower bounds are attained in the proximity
partial-sum order.  In the
$p=3m+1$ and $p=3$ cases, equality at the terminal raw coefficient forces
all $H_q=h$ and every branch coefficient to be $m$ and $1$, respectively.
In the $p=3m+2$ case, attaining the lower bound in (3.4) likewise forces
$H_q=h$ for all $q$; otherwise the terminal strict coefficient would fall
below $p^3$.  The branch sum must then be at least
\[
 p^3-A-(p-1)h=Lm+pm.
\]
Among nonincreasing integer sequences of length $L$, bounded above by
$m+1$, and with sum at least $Lm+pm$, the sequence with $pm$ entries
$m+1$ followed by $L-pm$ entries $m$ has the smallest prefix sums.
It is therefore the least sequence in the order imposed by proximity.
Substitution in
$A+\sum_{q\leq j}H_q$ and in the branch prefix sums verifies every
intermediate raw bound
$p(j+1)$ and $p(p+r)$, with equality at the endpoint.  Hence no coefficient
of any antinef majorant can be lowered, proving (3.2).

The Hoskin--Deligne formula \cite{LIP93,LW03} now gives
\[
 M_p=\binom{A+1}{2}+(p+1)(p-1)\binom{H+1}{2}
 +(p+1)p\sum_{i=1}^{L}\binom{K_i+1}{2}.
\tag{3.5}
\]
Substitution and simplification yield (3.3).  For $p=3$, this is
$120+48+72=240$; no finite profile is used in the derivation.
\end{proof}

\begin{corollary}[First transfer congruence]
\label[corollary]{cor:first-transfer-congruence}
For every prime $p$,
\[
 \tau_3(\A_2)\subseteq\LL_3+p\A_3.
\]
\end{corollary}

\begin{proof}
Under the coefficient pairing, the displayed inclusion is equivalent to
$J_3\subseteq(x^{p^2},y^{p^2})$.  For $p\geq5$, the augmentation order of the complete
ideal $J_3$ is the root point-basis coefficient $A$ in (3.2), and those
formulas give $A\geq2p^2-1$ for $p\geq5$.  A monomial whose two exponents
are both below $p^2$ has total degree at most $2p^2-2$; hence every element
of $J_3$ lies in $(x^{p^2},y^{p^2})$.  For $p=2,3$, the two-lane exact
certificate in \cref{app:exceptional-branches} constructs one point-fibre
indicator in the modular depth-three line code and independently verifies
its orthogonality to the full dual nullspace.  Affine translation gives all
point fibres, which is the same inclusion.  Thus the congruence holds for
every prime.
\end{proof}

The first cyclic diagonal block in the relative top shell has a unit count
$a_0$ which must not be confused with $m_0=M_p$.

\begin{proposition}[Cyclic unit count]
\label[proposition]{prop:a0}
\[
 a_0=\begin{cases}
 8,&p=2,\\
 108,&p=3,\\[1mm]
 \dfrac{(p-1)(3p^5+2p^4+2p^2+2)}{18},
 &p>3,\ p\equiv1\pmod3,\\[3mm]
 \dfrac{p(p+1)(12p^4-16p^3+10p^2-3p+1)}{72},
 &p\equiv2\pmod3.
 \end{cases}
\tag{3.6}
\]
\end{proposition}

\begin{proof}
Write a cyclic index as $j=a+bp+cp^2$.  The Weber--K\"unzer valuation
\cite{WK04} becomes
\[
 \beta_p(j)=a+(2p-1)b+p(3p-2)c.
\]
The unit count is the sum of $p^2(p-1)-\beta_p(j)$ over the region where
this quantity is positive.  For $p=6h+1$, the region consists of
$0\leq c\leq2h-1$ with all $a,b$, followed by the two boundary slices
$c=2h$ and $b=2h$.  For $p=6h+5$, it consists of
$0\leq c\leq2h$ with all $a,b$ and the slice
$c=2h+1$, $0\leq b\leq5h+3$.  Summing these arithmetic progressions gives
(3.6); $p=2,3$ are the degenerate boundary cases and are evaluated
directly.
\end{proof}

\section{Tripotent sectors and the reduced \texorpdfstring{$q$}{q}-Pascal
\texorpdfstring{\protect\\}{ }common image}
\label{sec:sector}

Fix a prime $p$.  Put
\[
 k=\F_p,\qquad N=p^3,\qquad N_0=p^2,\qquad e=N-N_0,
 \qquad \overline\OO=k[s]/(s^e),\qquad\zeta=1+s.
\]
The phrase ``$q$-Pascal'' is a conventional name.  Whenever a Gaussian
polynomial variable is required below it is denoted by $\mathfrak q$; the
letters $r,t,u$ denote residue-class labels and $h$ denotes an ordinary
integer in a congruence decomposition of $p$.  Thus no residue integer is
also the $q$-Pascal indeterminate.

Define the source and the two different chart targets by
\[
 \begin{aligned}
 U_N^X&=k[X]/((X-1)^e),&
 R_N^Y&=k[Y]/((Y-1)^N),\\
 U_N^Y&=k[Y]/((Y-1)^e),&
 D_1&=U_N^X\otimes_kR_N^Y,\\
 E_{\mathcal A}&=\overline\OO^{\Z/N\Z},&
 E_{\mathcal X}&=\overline\OO^{\Z/N_0\Z}.
 \end{aligned}
\]
The modular chart maps have the explicit types
\[
 \begin{aligned}
 \mathcal A:D_1&\longrightarrow E_{\mathcal A},&
 \mathcal A(b)_t&=b(\zeta,\zeta^t),\quad t\in\Z/N\Z,\\
 \mathcal X:D_1&\longrightarrow E_{\mathcal X},&
 \mathcal X(b)_u&=b(\zeta^{pu},\zeta),\quad u\in\Z/N_0\Z.
 \end{aligned}
\tag{4.1}
\]
Let $\rho_Y:R_N^Y\twoheadrightarrow U_N^Y$ be truncation and let
$\iota_Y:U_N^Y\to R_N^Y$ be the degree-below-$e$ $k$-linear section.  Put
\[
 \pi=\id\otimes(\iota_Y\rho_Y),\qquad
 V_0=\ker\pi,\qquad E=\im\pi=U_N^X\otimes_kU_N^Y.
\]
On $E$ let $\tau$ exchange the two truncated tensor factors, and extend it
to the tripotent $S=\tau\pi:D_1\to D_1$.  Then $S^2=\pi$, $S^3=S$,
$S|_{V_0}=0$, and $S|_E=\tau$.  Direct evaluation of a multiple of
$(Y-1)^e$ shows
\[
 \mathcal A(V_0)=\mathcal X(V_0)=0,\qquad
 \mathcal A=\mathcal A\pi,\qquad \mathcal X=\mathcal X\pi.
\]
Write $K_{\mathcal A}=\ker\mathcal A$ and
$K_{\mathcal A\mathcal X}=\ker\mathcal A\cap\ker\mathcal X$.

\begin{definition}[First integral obstruction]
\label[definition]{def:epsilon-one}
The depth-three obstruction used in the Smith assembly is
\[
 \eps_1(p,3):=\rank_k
 \begin{bmatrix}\mathcal A\\ \mathcal A S\end{bmatrix}
 -\rank_k\mathcal A
 =\rank_k\bigl((\mathcal AS)|_{K_{\mathcal A}}\bigr).
\tag{4.2}
\]
The second equality follows by projecting the stacked image onto its first
coordinate.
\end{definition}

Let
\[
 \jmath_{\mathrm{nu}}:E_{\mathcal X}\hookrightarrow E_{\mathcal A},
 \qquad
 (\jmath_{\mathrm{nu}}f)_t=
 \begin{cases}f_u,&t=pu\text{ with }u\in\Z/N_0\Z,\\
 0,&p\nmid t.
 \end{cases}
\tag{4.3}
\]
Every nonunit class $t\bmod N$ has a unique expression $t=pu$ with
$u\bmod N_0$, so this is a well-defined injective $k$-linear map.  The
notation $\jmath_{\mathrm{nu}}^{-1}$ below means its inverse only on the
supported subspace $\im\jmath_{\mathrm{nu}}$.

From this point through \cref{sec:newton} assume that $p$ is odd.  The
definition in \cref{def:epsilon-one} and the characteristic-free tripotent
rank identity remain valid at $p=2$, but the eigenspace splitting and the
factor $1/2$ below do not.

\begin{theorem}[Common-image reduction]
\label[theorem]{thm:common-image}
Let $E=\im\pi=E_+\oplus E_-$ be the $\pm1$ eigenspace decomposition of
$\tau$, and put
\[
 \mathcal A_+=\mathcal A|_{E_+},\qquad
 \mathcal A_-=\mathcal A|_{E_-},\qquad
 J_p^{\mathrm{full}}=\im\mathcal A_+\cap\im\mathcal A_-.
\]
These images and their intersection are $k$-subspaces of
$E_{\mathcal A}$; no $\overline\OO$-module structure is used here.  Define
\[
 \begin{aligned}
 \Fib_p&=\ker[\mathcal A_+,-\mathcal A_-]
       \subset E_+\oplus E_-,\\
 N_{\Fib}&=\ker\mathcal A_+\oplus\ker\mathcal A_-,\\
 c_{\Fib}:\Fib_p/N_{\Fib}&\longrightarrow J_p^{\mathrm{full}},&
 [(x,y)]&\longmapsto\mathcal A_+x=\mathcal A_-y.
 \end{aligned}
\]
Then $c_{\Fib}$ is an isomorphism and
\[
 \eps_1(p,3)=\dim_kJ_p^{\mathrm{full}}.
\tag{4.4}
\]
Moreover, $J_p^{\mathrm{full}}$ is supported on the nonunit slopes and the
transport $2\jmath_{\mathrm{nu}}^{-1}$ identifies it with
\[
 J_p^{\mathrm{red}}:=
 2\jmath_{\mathrm{nu}}^{-1}(J_p^{\mathrm{full}})
 =\mathcal X(K_{\mathcal A})\subset E_{\mathcal X}.
\tag{4.5}
\]
\end{theorem}

\begin{proof}
For any tripotent $S$, projection of
$v\mapsto(\mathcal Av,\mathcal ASv)$ onto its first coordinate shows that
the rank gap is $\rank(\mathcal AS|_{\ker\mathcal A})$.  Since
$\mathcal A(V_0)=0$, the zero sector disappears.  On $E_+\oplus E_-$ the
invertible target change $(a,b)\mapsto((a+b)/2,(a-b)/2)$ turns the stacked
map into the direct sum of the two sector maps.  Subtracting the rank of
$\mathcal A$ leaves precisely the dimension of their common image.  The
map $c_{\Fib}$ is onto by the definition of the intersection, and its
kernel before quotienting is exactly $N_{\Fib}$.

Put $P_+=(\pi+S)/2$ and $P_-=(\pi-S)/2$.  The map
\[
 \Gamma:K_{\mathcal A}\longrightarrow\Fib_p,\qquad
 b\longmapsto(P_+b,-P_-b)
\]
is onto: a preimage of $(x,y)\in\Fib_p$ is $b=x-y$.  It has kernel $V_0$,
and
\[
 c_{\Fib}[\Gamma(b)]=\frac12\mathcal ASb.
\tag{4.6}
\]
For $b\in K_{\mathcal A}$, the exact exchange identity is
\[
 (\mathcal ASb)_t=0\quad(p\nmid t),\qquad
 (\mathcal ASb)_{pu}=\mathcal X(b)_u.
\]
Consequently
\[
 c_{\Fib}\Gamma=\frac12\jmath_{\mathrm{nu}}
 \mathcal X|_{K_{\mathcal A}}.
\tag{4.7}
\]
Surjectivity of $c_{\Fib}\Gamma$ proves the support assertion and (4.5).
The factor $2$ is a unit because $p$ is odd.
\end{proof}

The reduced image has a sparse $q$-Pascal presentation.  For
$0\leq i<N$ and $0\leq k<e$, define
\[
 E_p(i,k)_r=(1-\zeta^{pr})^kC_i(r),\qquad
 C_i(r)=\sum_{j=i}^{N-1}\zeta^j
 \genfrac{[}{]}{0pt}{}{j}{i}_{\zeta^{-pr}},
 \quad 0\leq r<p^2.
\tag{4.8}
\]
Let
\[
 \beta_p(n)=\sum_{d=1}^{n}p^{v_p(d)},\qquad
 \lambda_i=\beta_p(N-1-i).
\tag{4.9}
\]
The source labels satisfy $\max(e-\lambda_i,0)\leq k<e$.

\begin{corollary}[$q$-Pascal presentation]
\label[corollary]{cor:qpascal}
After removing only columns proved identically zero,
\[
 J_p^{\mathrm{red}}=\im(E_p),\qquad
 \eps_1(p,3)=\rank_k(E_p).
\tag{4.10}
\]
\end{corollary}

\begin{proof}
The reduced fibre product $\Fib_p/N_{\Fib}$ is $J_p^{\mathrm{full}}$ by
\cref{thm:common-image}.  The Newton basis of $K_{\mathcal A}$ supplies the
columns (4.8), and the exchange identity in
the proof of \cref{thm:common-image} identifies their image with the
transported fibre product.  Quotienting by the column-relation space proves
(4.10).
\end{proof}

\section{Exact \texorpdfstring{$\mathfrak q$}{q}-Lucas block product
\texorpdfstring{\protect\\}{ }and support}
\label{sec:block}

For $r\in\Z/p^2\Z$, put $Q_r=\zeta^{-pr}$ and let
\[
 L=L_r=\ord(Q_r)=
 \begin{cases}1,&r=0,\\p,&r\neq0,\ p\mid r,\\p^2,&p\nmid r,
 \end{cases}
 \qquad M=N/L.
\]
Write $i=dL+c$ with $0\leq c<L$.

\begin{lemma}[$\mathfrak q$-Lucas evaluation]
\label[lemma]{lem:qlucas}
The evaluation $\mathfrak q\mapsto Q_r$ kills $p$ and
$\Phi_L(\mathfrak q)$.  Hence, for
$0\leq b<L$,
\[
 \genfrac{[}{]}{0pt}{}{aL+b}{dL+c}_{Q_r}
 =\binom ad\genfrac{[}{]}{0pt}{}{b}{c}_{Q_r}.
\tag{5.1}
\]
\end{lemma}

\begin{proof}
If $L=p^a$, then in characteristic $p$,
$\Phi_{p^a}(Z)=(Z-1)^{p^{a-1}(p-1)}$.  For $L=p^2$ the order of
$Q_r-1$ is $p$, while for $L=p$ it is $p^2$; in both cases the displayed
power has order $e$ and vanishes in $\overline\OO$.  Sagan's congruence
\cite[Theorem~2.2]{SAG92} therefore applies.  The case $L=1$ is ordinary
Lucas.
\end{proof}

Two finite identities will be used:
\begin{align}
 \sum_{a=d}^{M-1}\binom adx^a
 &=x^d(1-x)^{M-1-d},\tag{5.2}\label{eq:ordinary-block}\\
 \sum_{b=c}^{L-1}x^b\genfrac{[}{]}{0pt}{}{b}{c}_{Q}
 &=x^c\prod_{h=c+1}^{L-1}(1-xQ^h),\tag{5.3}\label{eq:cyclotomic-block}
\end{align}
where $M$ is a power of $p$ and $\Phi_L(Q)=0$.  The first follows by
extracting a coefficient from a finite geometric series and using
Frobenius.  The second is the finite $\mathfrak q$-binomial theorem after the standard
root-of-unity symmetry
$ \genfrac{[}{]}{0pt}{}{L-1-c}{t}_{Q}
=(-1)^tQ^{-ct-t(t+1)/2}\genfrac{[}{]}{0pt}{}{c+t}{c}_{Q}$.

\begin{theorem}[Exact block product and support]
\label[theorem]{thm:block-product}
For every odd prime, slope and index,
\[
 C_i(r)=\zeta^i(1-\zeta^L)^{M-1-d}
 \prod_{h=c+1}^{L-1}(1-\zeta^{1-prh}).
\tag{5.4}
\]
Consequently
\[
 C_i(r)=
 \begin{cases}
 (-1)^is^{N-1-i}+O(s^{N-i}),&N-1-i<e,\\
 0,&N-1-i\geq e,
 \end{cases}
\tag{5.5}
\]
 and the column $(i,k)$ in (4.8) is nonzero exactly when
\[
 \max(e-\lambda_i,0)\leq k<e,\qquad N-1-i+pk<e.
\tag{5.6}
\]
\end{theorem}

\begin{proof}
Write $j=aL+b$ in (4.8).  Applying \cref{lem:qlucas} separates the sum
into the two factors in \eqref{eq:ordinary-block} and
\eqref{eq:cyclotomic-block}, which gives (5.4).  Since
$1-\zeta^L=-s^L$ and every $1-\zeta^{1-prh}$ begins with $-s$, its total
order is
\[
 L(M-1-d)+(L-1-c)=N-1-i.
\]
The parity of the number of negative factors is the parity of $i$, proving
(5.5).  On a unit slope,
$1-\zeta^{pr}=-r_0s^p+O(s^{2p})$; this proves sufficiency in (5.6).
Nonunit slopes have no smaller prefactor order, proving necessity.
\end{proof}

The equality (5.4) belongs to the truncated target $\overline\OO$; it is
not an equality in the unquotiented group algebra.  It resolves cancellation
inside a single column but does not yet control relations among columns.

\section{Boundary generation and the
\texorpdfstring{\protect\\}{ }Newton-minors theorem}
\label{sec:newton}

Put $u=\zeta^p=1+s^p$ and $T_r=1-u^r$.  Reverse the column index by
\[
 \mathcal C_t(r)=C_{N-1-t}(r),\qquad
 F_{t,K}(r)=T_r^K\mathcal C_t(r).
\tag{6.1}
\]
The surviving labels are exactly
\[
 K+\beta_p(t)\geq e,\qquad t+pK<e.
\tag{6.2}
\]
Let $W=\im(E_p)\subset\overline\OO^{p^2}$, initially only as a $k$-space.

\begin{lemma}[Scalar recurrence and module closure]
\label[lemma]{lem:recurrence}
With $c=s\zeta^{-1}$,
\[
 cF_{t,K}=\sum_{\ell=1}^{t+1}(-1)^{\ell+1}
 \binom{t+1}{\ell}F_{t,K+\ell}-F_{t+1,K}.
\tag{6.3}
\]
All coefficients are in $k$ and independent of $r$.  Consequently $W$ is
an $\overline\OO$-submodule of $\overline\OO^{p^2}$.
\end{lemma}

\begin{proof}
The exact recurrence before expansion is
$F_{t+1,K}=-cF_{t,K}+[t+1]_{u^r}F_{t,K+1}$.  Since
\[
 T_r[t+1]_{u^r}=1-(1-T_r)^{t+1},
\]
binomial expansion gives (6.3).  Every nonzero term remains admitted or is
zero by the support boundary, so $cW\subseteq W$.  Finally $c^e=0$ and
$s=c/(1-c)=c+\cdots+c^{e-1}$, hence $sW\subseteq W$.
\end{proof}

Define
\[
 \mu_p(x)=\min\{t\geq0:\beta_p(t)\geq x\},\qquad
 t_K=\mu_p(e-K),\qquad G_K=F_{t_K,K},
\tag{6.4}
\]
and
\[
 a_K=t_K+pK,\qquad
 H_p=\#\{0\leq K<p(p-1):a_K<e\}.
\tag{6.5}
\]
Here $G_K$ is a vector and $H_p$ is an integer.  These types are fixed
throughout the paper.

\begin{lemma}[Boundary generation]
\label[lemma]{lem:boundary}
The nonzero boundary indices are $0\leq K<H_p$, and
\[
 W=\sum_{K=0}^{H_p-1}\overline\OO G_K.
\tag{6.6}
\]
\end{lemma}

\begin{proof}
Since $\mu_p(x)-\mu_p(x-1)\in\{0,1\}$,
$a_{K+1}-a_K\geq p-1$, so the nonzero indices form an initial interval.
For the induction in $t$, one needs the following Lucas admission step:
if $\binom t\ell\not\equiv0\pmod p$, then
\[
 \ell\geq p^{v_p(t)}=\beta_p(t)-\beta_p(t-1).
\]
Indeed, Lucas forces the first $v_p(t)$ digits of $\ell$ to vanish.  Apply
(6.3) at $(t-1,K)$ and solve for $F_{t,K}$.  Every term with nonzero
coefficient remains admitted by the displayed inequality and has smaller
first index; induction reaches $G_K$.
\end{proof}

Lift to $R=k[[s]]$.  Define
\[
 \widetilde{\mathcal C}_t(r)=\zeta^{-1}
 \prod_{\ell=1}^{t}(\zeta^{-1}-u^{r\ell}),\qquad
 \widetilde G_K(r)=T_r^K\widetilde{\mathcal C}_{t_K}(r).
\tag{6.7}
\]
Put $X_r=[r]_u$ and
$Q_j(r)=\genfrac{[}{]}{0pt}{}{r}{j}_u$.  The matrix
$(Q_j(r))_{0\leq r,j<p^2}$ is lower unitriangular, so the $Q_j$ form an
$R$-basis of $R^{p^2}$.

\begin{lemma}[Coefficient cost]
\label[lemma]{lem:cost}
There are a unit $\eta_t(s)$ and coefficients $h_{t,m}(s)$ such that
\[
 \widetilde{\mathcal C}_t
 =s^t\eta_t(s)\sum_{m\geq0}h_{t,m}(s)X^m,\qquad
 h_{t,0}=1,\qquad v_s(h_{t,m})\geq(p-1)m.
\tag{6.8}
\]
\end{lemma}

\begin{proof}
Each factor in (6.7) is
\[
 -c+\sum_{d=1}^{\ell}(-1)^{d+1}\binom\ell d(1-u)^dX^d.
\]
After the constant term is factored, a selection with total $X$-degree
$m$ has extra $s$-order at least $pm$ minus the number of selected factors,
hence at least $(p-1)m$.  The all-constant selection is unique.
\end{proof}

Carlitz's integral identity \cite{CER18} is
\[
 X^n=\sum_{j=0}^{n}u^{\binom j2}S_u[n,j][j]_u!\,Q_j.
\tag{6.9}
\]
For $j<p^2$, put
\[
 \phi_j=v_s([j]_u!)=p(\beta_p(j)-j),\qquad
 \tau_p(K)=a_K+\phi_K=\mu_p(e-K)+p\beta_p(K).
\tag{6.10}
\]
The sequence $\tau_p(K)$ is strictly increasing.

\begin{theorem}[Newton-minors theorem]
\label[theorem]{thm:newton-minors}
Let $\mathcal M_p$ be the $p^2\times H_p$ matrix of the lifted boundary
columns in the $\mathfrak q$-Newton row basis.  If $I_r$ is its ideal of $r\times r$
minors, then
\[
 v_s(I_r)=\sum_{K=0}^{r-1}\tau_p(K)\qquad(1\leq r\leq H_p).
\tag{6.11}
\]
Thus the Smith exponents over the DVR $R$ are exactly
$\tau_p(0),\ldots,\tau_p(H_p-1)$.
\end{theorem}

\begin{proof}
Combining \cref{lem:cost} with (6.9), every path
$K\to K+m\to j$ contributing to row $j$ of column $K$ has order at least
$a_K+\phi_j$; equality occurs for $m=0,j=K$.  For column indices
$K_1<\cdots<K_r$ and row indices $J_1<\cdots<J_r$, every determinant
monomial therefore has order at least
\[
 \sum_i a_{K_i}+\sum_i\phi_{J_i}
 \geq\sum_{K=0}^{r-1}\tau_p(K).
\]
In the principal minor on rows and columns $0,\ldots,r-1$, the identity
permutation with paths $m=0,j=K$ attains this bound uniquely.  A path with
$m>0$ has positive extra cost; with $m=0$, Carlitz support requires
$\pi(K)\leq K$ for each column, forcing the identity permutation.  Hence
the minimum term cannot cancel, proving (6.11).  Successive differences of
determinantal-ideal valuations are the DVR Smith exponents.
\end{proof}

\section{Arithmetic evaluation and the filtered common image}
\label{sec:filtered}

Reduction of the DVR Smith form modulo $s^e$, together with
\cref{lem:boundary}, gives the obstruction dimension
\[
 \eps_1(p,3)=\sum_{K=0}^{p(p-1)-1}
 \pos{e-\mu_p(e-K)-p\beta_p(K)}.
\tag{7.1}
\]

Put
\[
 L_K=e-\tau_p(K)=e-\mu_p(e-K)-p\beta_p(K).
\tag{7.2}
\]

\begin{lemma}[Exact inversion of $\beta_p$]
\label[lemma]{lem:beta-inversion}
Let $p>3$, write $p=6h+\rho$ with $\rho\in\{1,5\}$, and set
$d=2p-1$.  For $0\leq K\leq M$ choose the branch constants
\[
\begin{array}{c|c|c|c|c}
\rho&A&B_0&c&M\\ \hline
1&2h&2h&4h&h(2p+1)\\[1mm]
5&2h+1&5h+4&7h+5&\dfrac{p^2-1}{3}
\end{array}
\tag{7.3}
\]
and define
\[
 \begin{aligned}
 R_K&=e-K-Ap(3p-2),\\
 B_K&=\left\lceil\frac{R_K-(p-1)}d\right\rceil
     =B_0-\left\lfloor\frac{K+c}{d}\right\rfloor,\\
 C_K&=\max(0,R_K-dB_K),\\
 \delta_K&=\max(0,dB_K-R_K)
 =\max\bigl(0,((K+c)\bmod d)-(p-1)\bigr).
 \end{aligned}
\tag{7.4}
\]
Then $A,B_K,C_K$ are the base-$p$ digits of the least integer whose
$\beta_p$-value is at least $e-K$, and
\[
 \mu_p(e-K)=Ap^2+R_K-(p-1)B_K+\delta_K.
\tag{7.5}
\]
Consequently
\[
 \boxed{L_K=(p-1)\left(2Ap+B_0-K
 -p\left\lfloor\frac Kp\right\rfloor
 -\left\lfloor\frac{K+c}{2p-1}\right\rfloor\right)-\delta_K.}
\tag{7.6}
\]
\end{lemma}

\begin{proof}
For $0\leq t<p^3$ the valuation-count identity is
\[
 \beta_p(t)=t+(p-1)\left\lfloor\frac tp\right\rfloor
 +(p^2-p)\left\lfloor\frac t{p^2}\right\rfloor.
\]
Writing $t=Ap^2+Bp+C$ with $0\leq A,B,C<p$ therefore gives
\[
 \beta_p(t)=Ap(3p-2)+B(2p-1)+C.
\tag{7.7}
\]
For the intervals in (7.3), direct substitution at $K=0,M$ keeps the
first digit equal to the displayed $A$.  Once $A$ is fixed, the least
$B$ for which the residual target can be reached with $0\leq C<p$ is
$B_K$ in (7.4), and the least last digit is $C_K$.  Since
$C_K=R_K-dB_K+\delta_K$, substitution into $Ap^2+B_Kp+C_K$ gives (7.5).
The identity $R_0-(p-1)=B_0d-c$ gives the second formula for $B_K$ and the
residue formula for $\delta_K$.  Finally, for $K<p^2$,
$\beta_p(K)=K+(p-1)\lfloor K/p\rfloor$; inserting this and (7.5) into
(7.2) gives (7.6).
\end{proof}

\begin{lemma}[Positive-tail endpoints and residue sums]
\label[lemma]{lem:tail-endpoints}
With the branch constants of \cref{lem:beta-inversion},
\[
 L_{M-1}=p-1,\qquad L_M=0,\qquad
 L_K>0\Longleftrightarrow 0\leq K<M.
\tag{7.8}
\]
Moreover, the two floor sums and the periodic ramp in (7.6) are
\[
\begin{array}{c|c|c}
&p=6h+1&p=6h+5\\ \hline
\displaystyle\sum_{K=0}^{M-1}\left\lfloor\frac Kp\right\rfloor
&ph(2h-1)+2h^2
&ph(2h+1)+(2h+1)(4h+3)\\[3mm]
\displaystyle\sum_{K=0}^{M-1}\left\lfloor\frac{K+c}{2p-1}\right\rfloor
&(2p-1)\dfrac{h(h-1)}2+6h^2
&(2p-1)\dfrac{h(h+1)}2+(h+1)(p-1)\\[3mm]
\displaystyle\sum_{K=0}^{M-1}\delta_K
&3h^2p
&(h+1)\dfrac{p(p-1)-h}{2}.
\end{array}
\tag{7.9}
\]
\end{lemma}

\begin{proof}
Substitution of (7.3) into (7.6) gives the two endpoint equalities.
Strict increase of $\tau_p(K)$ from \cref{thm:newton-minors} makes $L_K$
strictly decreasing, so they prove the last assertion of (7.8).  For the
first row of (7.9), write $M=ap+b$ and sum $a$ complete residue blocks.
For the second and third rows, write $K+c=j(2p-1)+r$ and sum separately
over $0\leq r<2p-1$ and the terminal block; the ramp is nonzero precisely
for $p\leq r<2p-1$.  The complete residue decompositions are displayed in
\cref{app:arithmetic-blocks}; they give (7.9) without an asymptotic or
interpolation step.
\end{proof}

\begin{theorem}[Closed odd-prime obstruction]
\label[theorem]{thm:epsilon}
For odd primes,
\[
 \eps_1(p,3)=\FCIGDThreeEpsilonClosedFormula
\tag{7.10}
\]
\end{theorem}

\begin{proof}
For $p>3$, \cref{lem:tail-endpoints} reduces (7.1) to
$\sum_{K=0}^{M-1}L_K$.  Summing (7.6), using
$\sum_{K=0}^{M-1}K=M(M-1)/2$, and substituting the three exact rows of
(7.9) gives the two polynomials in (7.10).  For $p=3$, the active
thresholds are $9,12,15$ and $e=18$, so the sum is
$(18-9)+(18-12)+(18-15)=18$.
\end{proof}

\begin{theorem}[Filtered common-image profile]
\label[theorem]{thm:filtered-profile}
As an $\overline\OO$-module,
\[
 W\cong\bigoplus_{\substack{0\leq K<H_p\\\tau_p(K)<e}}
 \overline\OO/(s^{e-\tau_p(K)}).
\tag{7.11}
\]
For the ambient filtration
$F^d_{\mathrm{amb}}W=W\cap s^d\overline\OO^{p^2}$,
\[
 \dim_k(F^d_{\mathrm{amb}}W/F^{d+1}_{\mathrm{amb}}W)
 =\#\{0\leq K<H_p:\tau_p(K)\leq d<e\}.
\tag{7.12}
\]
Hence
\[
 \Hilb_W^{\mathrm{amb}}(z)=
 \sum_{\substack{0\leq K<H_p\\\tau_p(K)<e}}
 z^{\tau_p(K)}\frac{1-z^{e-\tau_p(K)}}{1-z}.
\tag{7.13}
\]
The minimal number of generators is
\[
 \ngen_{\overline\OO}(W)=\dim_k(W/sW)=
 \begin{cases}
 3,&p=3,\\[1mm]
 \dfrac{(p-1)(2p+1)}6,&p\equiv1\pmod6,\\[3mm]
 \dfrac{p^2-1}{3},&p\equiv5\pmod6.
 \end{cases}
\tag{7.14}
\]
\end{theorem}

\begin{proof}
Invertible Smith row operations over $R$ preserve every ambient power
$s^dR^{p^2}$.  Reduction modulo $s^e$ turns the active diagonal with
threshold $\tau_p(K)$ into the embedded ideal
$s^{\tau_p(K)}\overline\OO$, proving (7.11)--(7.12).  Summing the successive
dimensions gives (7.13).  Nakayama's lemma counts one generator for each
nonzero cyclic summand, and the positive-tail counts in the proof of
\cref{thm:epsilon} give (7.14).
\end{proof}

The boundary-column count and generator count are different invariants:
\[
\begin{array}{c|rrrr}
p&3&5&7&11\\ \hline
H_p&3&12&28&71\\
\ngen_{\overline\OO}(W)&3&8&15&40.
\end{array}
\tag{7.15}
\]

\begin{corollary}[Truncation redundancy]
\label[corollary]{cor:redundancy}
\[
 H_p-\ngen_{\overline\OO}(W)
 =\#\{0\leq K<H_p:\tau_p(K)\geq e\}
\tag{7.16}
\]
counts the Smith directions in the boundary presentation eliminated by
truncation modulo $s^e$.
\end{corollary}

This is a numerical structural statement.  It does not assert that a
particular named subset of the original boundary columns is zero.

\section{The cross-chart transfer defect}
\label{sec:o2}

Let $\eta_2:\QQ_2/p\QQ_2\to p\QQ_3/p^2\QQ_3$ be induced by transfer and
put $d_p=\dim_k\ker\eta_2$.  The equality
\[
 \tau_3(\QQ_2)\cap p^2\QQ_3=\tau_3(p\QQ_2)
\tag{O2}
\]
is equivalent to $d_p=0$.

\begin{lemma}[Defect bookkeeping]
\label[lemma]{lem:defect}
Let $T_{\geq3}=m_3+m_4+m_5$.  Then
\[
 T_{\geq3}=a_0+\eps_1(p,3)+d_p.
\tag{8.1}
\]
\end{lemma}

\begin{proof}
The multiscale identity gives $p^3\QQ_3=\tau_3(p\QQ_2)$.  The natural map
$p^2\QQ_3/p^3\QQ_3\to p^2(\QQ_3/\tau_3\QQ_2)$ is onto, with kernel
\[
 \frac{\tau_3\QQ_2\cap p^2\QQ_3}{\tau_3(p\QQ_2)},
\]
whose dimension is $d_p$.  The relative top-shell calculation gives
$a_0+\eps_1$ for the target dimension, proving (8.1).
\end{proof}

We now define the carry objects used to kill this defect.  Let
\[
 \mathscr R_2=k[G,H]/(G^{N_0}-1,H^{N_0}-1),\qquad
 x=G-1,\quad y=H-1,
\]
identified with the $k$-space of functions on the coarse point set
$G_2=(\Z/N_0\Z)^2$.  With $\iota([g])=[-g]$, use the Frobenius pairing
$\langle a,b\rangle=[1](a\iota(b))$.  For $\delta\in\mathbb P^1(k)$ choose
\[
 \begin{array}{c|cc}
 \delta&g_\delta&t_\delta\\ \hline
 \delta\in k&GH^\delta&H\\
 \infty&H&G,
 \end{array}
 \qquad z_\delta=g_\delta-1,\quad v_\delta=t_\delta-1,
\]
and define the tangent-bundle ideal and line code
\[
 I_\delta=(z_\delta,v_\delta^p)^p,\qquad
 \mathcal C_{2,\delta}=I_\delta^\perp,\qquad
 \mathcal C_2=\sum_{\delta\in\mathbb P^1(k)}\mathcal C_{2,\delta}.
\tag{8.2}
\]
The normal-to-tangent conversion is
\[
 \kappa(0)=\infty,\qquad \kappa(d)=-d^{-1}\ (d\in k^\times),
 \qquad \kappa(\mathrm{cross})=0.
\tag{8.3}
\]

For $b\in K_{\mathcal A}$ choose a compatible integral lift
$\widetilde b$ in the fixed cyclotomic lattice and write its first-chart
remainders as
\[
 A_t(\widetilde b)=\sum_{j<e}a_{t,j}\zeta^j.
\]
Every $a_{t,j}$ is divisible by $p$.  If moreover
$b\in K_{\mathcal A\mathcal X}$, write
\[
 X_u(\widetilde b)=\sum_{j<e}x_{u,j}\zeta^j;
\]
then every $x_{u,j}$ is divisible by $p$ as well.  If $o\in\Z/N\Z$, let
$[c_{\bullet}]_{o}^{<e}$ mean $c_{[o]_N}$ when $0\leq[o]_N<e$ and zero
when $e\leq[o]_N<N$.  For $(r,s)\in G_2$ put
\[
 o_t(r,s)=[e+r+t(e+s)]_N,\qquad
 o_u^{\mathcal X}(r,s)=[pu(e+r)+e+s]_N.
\tag{8.4}
\]
The first-chart residue carries and the cross carry are the typed maps
\[
 \begin{gathered}
 \mathcal H_d:K_{\mathcal A}\longrightarrow\mathscr R_2
 \quad(d\in k),\\
 \mathcal H_X:K_{\mathcal A\mathcal X}\longrightarrow\mathscr R_2,\\
 \mathcal H_A:=\sum_{d\in k}\mathcal H_d:
 K_{\mathcal A}\longrightarrow\mathscr R_2.
 \end{gathered}
\]
Their values on compatible lifts are
\[
 \begin{aligned}
 \mathcal H_d(b)(r,s)
 &=\sum_{\substack{t\bmod N\\t\equiv d\bmod p}}
   \frac{[a_{t,\bullet}]_{o_t(r,s)}^{<e}}p\pmod p,\\
 \mathcal H_X(b)(r,s)
 &=\sum_{u\bmod N_0}
   \frac{[x_{u,\bullet}]_{o_u^{\mathcal X}(r,s)}^{<e}}p\pmod p.
 \end{aligned}
\tag{8.5}
\]
Division is performed only after an integral candidate relation has been
combined, so its coefficients are divisible by $p$.  If the compatible
lift is changed by $p h$, the change in (8.5) is an ordinary carrier in the
corresponding line bundle; hence the map
\[
 \delta_1:K_{\mathcal A\mathcal X}\longrightarrow
 \mathscr R_2/\mathcal C_2,
 \qquad
 \delta_1(b)=[\mathcal H_A(b)+\mathcal H_X(b)]
\tag{8.6}
\]
is independent of the lift.

For an integral representative define the one-way exchange
\[
 \widetilde S(\widetilde b)=
 \operatorname{Rem}_{\Phi_{p^3}(X)}\widetilde b(Y,X),
\tag{8.7}
\]
with monic remainder in $X$.  Modulo $p$,
$\Phi_{p^3}(X)=(X-1)^e$; hence this reduction is exactly the tripotent
$S=\tau\pi:D_1\to D_1$ defined in \cref{sec:sector}, rather than a second
operator with the same name.

\begin{lemma}[Carry transposition]
\label[lemma]{lem:carry-transpose}
For odd $p$,
\[
 S(K_{\mathcal A\mathcal X})\subset K_{\mathcal A},\qquad
 \mathcal H_0(Sb)(r,s)=\mathcal H_X(b)(s,r).
\tag{8.8}
\]
If $P$ is point transposition, $(Pf)(r,s)=f(s,r)$, then
\[
 P(\mathcal C_{2,\infty})=\mathcal C_{2,0}.
\tag{8.9}
\]
\end{lemma}

\begin{proof}
Monic division in (8.7) gives the integral identities
\[
 A_t(\widetilde S\widetilde b)=\widetilde b(\zeta^t,\zeta),
 \qquad A_{pu}(\widetilde S\widetilde b)=X_u(\widetilde b).
\tag{8.10}
\]
For a unit $t$, the automorphism $\zeta\mapsto\zeta^t$ transports the
vanishing of $\mathcal A(b)$ at $t^{-1}$ to the first expression in
(8.10); for a nonunit $t=pu$, the second kernel condition applies.  This
proves the inclusion.  The second identity in (8.10) identifies the
integer coefficient arrays, while
\[
 o_{pu}(r,s)=e+r+pu(e+s)=o_u^{\mathcal X}(s,r)\pmod N.
\]
Thus clipping, division and summation agree term by term, proving (8.8).
Finally, transposition takes the $p$ depth-two lines with normals $(1,pu)$
to those with normals $(pu,1)$, proving (8.9).
\end{proof}

To state the seed calculation, let
$\rho_p:D_1\to\mathscr R_2$ reduce both exponents modulo $N_0$, put
\[
 P_j(X,Y)=\prod_{a=0}^{j-1}(Y-X^a),\qquad
 b_{j,\nu}=(X-1)^\nu P_j(X,Y),\qquad
 \nu\geq\max(e-\beta_p(j),0),
\tag{8.11}
\]
and let $\Pi_0$ retain the monomials whose $Y$-degree is divisible by $p$.

\begin{lemma}[Newton kernel basis]
\label[lemma]{lem:newton-kernel-basis}
The vectors
\[
 b_{j,\nu},\qquad 0\leq j<N,\quad
 \max(e-\beta_p(j),0)\leq\nu<e,
\]
form a $k$-basis of $K_{\mathcal A}$.
\end{lemma}

\begin{proof}
Each $P_j$ is monic of degree $j$ in $Y$, so the transition from
$\{(X-1)^\nu Y^j\}$ to $\{(X-1)^\nu P_j\}$ is unitriangular over
$k[X]/((X-1)^e)$.  Moreover
\[
 P_j(\zeta,\zeta^t)
 =P_j(\zeta,\zeta^j)
   \genfrac{[}{]}{0pt}{}{t}{j}_{\zeta},
\]
where the Gaussian-Pascal matrix is lower unitriangular in $t,j$.
Invertible row operations therefore diagonalize first-chart evaluation with
diagonal entries
\[
 d_j=P_j(\zeta,\zeta^j)
 =\prod_{a=0}^{j-1}(\zeta^j-\zeta^a).
\]
For $1\leq d\leq j$, the factor $\zeta^d-1$ is a unit times
$s^{p^{v_p(d)}}$; hence $d_j$ is a unit times
$s^{\beta_p(j)}$.  In $\overline\OO=k[s]/(s^e)$ the kernel of
multiplication by $d_j$ is exactly
$s^{\max(e-\beta_p(j),0)}\overline\OO$, whose displayed powers of $s$
form a $k$-basis.  Taking the direct sum over $j$ proves the claim.
\end{proof}

For a
first-chart residue $d$, the $p+1$ mixed differences generating
$I_{\kappa(d)}$ give functionals
$\lambda_{d,c}(b)$, $0\leq c\leq p$, by applying
$z_{\kappa(d)}^{p-c}(v_{\kappa(d)}^p)^c$ to $\mathcal H_d(b)$.

\begin{lemma}[Seed-defect factorization and annihilation]
\label[lemma]{lem:seed-annihilation}
For every odd prime, every residue $d$ and $b\in K_{\mathcal A}$,
\[
 \lambda_{d,c}(b)=\rho_p(b)S_{d,c},\qquad0\leq c\leq p,
\tag{8.12}
\]
where, with $y_d=G^{-d}H-1$,
\begin{align}
 S_{d,0}&=-y_d^p+x^{p(p-1)}y_d^{p(p-1)}C_p(y_d),\notag\\
 S_{d,c}&=x^{pc}(y_d^{N_0-c}-y_d^{p-c})\quad(1\leq c<p),\notag\\
 S_{d,p}&=0,\qquad
 C_p(Y)=\sum_{a=1}^{p-1}\left(\frac1p\binom pa\bmod p\right)Y^a.
\tag{8.13}
\end{align}
The auxiliary carrier seeds are
\[
 \begin{aligned}
 B^0_{d,0}&=x^{p(p-1)}y_d^{p(p-1)}C_p(y_d),\\
 B^0_{d,c}&=x^{pc}y_d^{N_0-c}\quad(1\leq c<p),
 &B^0_{d,p}&=0.
 \end{aligned}
\]
They satisfy, on every admissible Newton basis vector,
\[
 \rho_p(b_{j,\nu})B^0_{d,c}=0,\qquad
 \rho_p(\Pi_0b_{j,\nu})B^0_{d,c}=0.
\tag{8.14}
\]
The second equality is the Cartier constant-term assertion.
\end{lemma}

\begin{proof}
Let $\sigma(n,o)$ be the coefficient of $\zeta^o$ in the standard
representative of $\zeta^n$ modulo $\Phi_{p^3}$, extended by zero for
$e\leq o<N$.  The last cyclotomic block gives exactly
\[
 \sigma(n,o)-\mathbf1_{o\equiv n\ (N)}
 =-\mathbf1_{[n]_N\geq e}\mathbf1_{o\equiv n\ (N_0)}.
 \tag{8.15}
\]
 We give the residue-packet calculation explicitly.  For a source monomial
 $m=X^iY^j$, put $J=[j]_p$ and, at the point $(r,s)$, put $S=[s]_p$.  Its
 undivided residue-$d$ carrier is
 \[
  T_d(m)(r,s)=\sum_{u\bmod N_0}
  \sigma\bigl(i+(d+pu)j,o_{d+pu}(r,s)\bigr).
 \tag{8.15a}
 \]
 Write $u=u_0+pv$, with $u_0,v\in k$.  The common $N_0$-tail condition is
 \[
  r+(d+pu_0)s\equiv i+(d+pu_0)j\pmod {N_0};
 \]
 denote its solution set in $u_0$ by $U$.  For $u_0\in U$, the source and
 clipped-point top-block digits have the form $A+Jv$ and $B+Sv$.  By
 (8.15), the complete $v$-packet is
 \[
  \mathsf P(A,B;J,S)=
  \sum_{v\in k}\mathbf1_{B+Sv\ne-1}
  \bigl(\mathbf1_{A+Jv=B+Sv}-\mathbf1_{A+Jv=-1}\bigr).
 \tag{8.15b}
 \]
 The four possibilities are as follows; the last column is evaluated for
 each $u_0\in U$.
 \[
 \begin{array}{c|c|c}
 (J,S)&|U|&\mathsf P(A,B;J,S)\pmod p\\ \hline
 J\ne S,\ J\ne0&0\text{ or }1&0\\
 J\ne S,\ J=0&0\text{ or }1&1\\
 J=S\ne0&0\text{ or }p&-1\quad(\sum_{u_0\in U}\mathsf P=0)\\
 J=S=0&0\text{ or }p&
 p\,\mathbf1_{B\ne-1}
 (\mathbf1_{A=B}-\mathbf1_{A=-1})=0\\
 \end{array}
 \]
 Indeed, in the first row the equations $A+Jv=B+Sv$ and $A+Jv=-1$
 each have one solution, and either both are excluded or neither is.  In the
 second row the first equation contributes once; if $A=-1$, its excluded
 contribution is replaced by $-(p-1)=1$.  In the third row the packet is
 $-1$ whether $A=B$ or not, and the $p$ values of $u_0$ cancel modulo $p$.
 The fourth row is constant in $v$ and is visibly a multiple of $p$.
 Consequently, if
 \[
  N_{d,u_0;i,j}=\{(r,s):
  r+(d+pu_0)s=i+(d+pu_0)j\pmod {N_0}\},
 \]
 then
 \[
  T_d(m)\bmod p=\mathbf1_{J=0}
  \sum_{u_0\in k}\mathbf1_{N_{d,u_0;i,j}}.
 \tag{8.15c}
 \]
 Thus the ordinary carrier is in the residue-$d$ line bundle, including all
 seam and wall contributions.

 To retain the information needed after division by $p$, let $L_d(m)$ be the
 integral $0$--$1$ line sum in (8.15c) and write
 $T_d=L_d+pR_d$.  With $\rho_p(m)=G^iH^j$, the same four rows before the
 final reduction give
 \[
  T_d(m)-\rho_p(m)T_d(1)\equiv
  \begin{cases}
  -\rho_p(m)L_d(1),&J\ne0,\\
  -p\displaystyle\sum_{u_0\in W_m}
       \mathbf1_{N_{d,u_0;i,j}},&J=0,
  \end{cases}
  \pmod {p^2},
 \tag{8.15d}
 \]
 where $W_m$ consists of the $u_0$ for which the constant source top-block
 digit is the last block.  For clarity, at the translated point
 $([r-i]_{N_0},[s-j]_{N_0})$ the seed has point digit $S-J$ and the same
 tail set $U$.  Canonicalizing that point changes only the constant term of
 the affine top-block test in $u_0$, not its linear coefficient.  If that
 coefficient is nonzero, both tests have one solution; if it is zero, each
 has zero or $p$ solutions.  Their counts therefore agree modulo $p$, and
 the complete $v$-packet supplies the second factor of $p$ in (8.15d).

 Let
 $D_{d,c}=z_{\kappa(d)}^{p-c}(v_{\kappa(d)}^p)^c$.  It annihilates every
 line sum in (8.15c) modulo $p$.  If $\mathcal B_{d,c}$ denotes the resulting
 line Bockstein, then on a combined first-chart relation
 \[
  \lambda_{d,c}=\mathcal B_{d,c}(L_d)+D_{d,c}R_d.
 \]
 Put $B^0_{d,c}=\mathcal B_{d,c}(L_d(1))$ and
 $S_{d,c}=\lambda_{d,c}(1)$.  Applying $D_{d,c}$ to (8.15d), dividing only
 after the complete relation has been combined, and extending linearly gives
 \[
  \lambda_{d,c}(b)-\rho_p(b)S_{d,c}
  =-\rho_p((1-\Pi_0)b)B^0_{d,c}.
 \tag{8.15e}
 \]
 The exact values of $B^0_{d,c}$ and $S_{d,c}$ are obtained by putting
$E=G^p$, $T=E-1=x^p$,
\[
 S_u(E)=\frac{E^u-1}{E-1},\qquad
 h(E)=\frac{(E-1)^p-(E^p-1)}p\pmod p.
\]
In $k[T]/(T^p)$ one has $h(1+T)=-\log(1+T)$.  For
$A_n(T)=\sum_{u\in k}E^uS_u(E)^n$, finite-field moments give
\[
 hA_n=0\ (0\leq n\leq p-2),\qquad hA_{p-1}=T,
\]
and, for $0\leq m\leq p-2$,
\[
 hA_{p+m}=\frac{(-1)^m}{m+1}T^{p-1-m}\pmod{T^{p-m}}.
 \tag{8.16}
\]
Indeed, with
$\phi(w)=\sum_{a=1}^{p-1}\binom waT^{a-1}$,
$\phi'(w)=-(h/T)(1+T)^w$ and $\phi(w)^p=w^p$ modulo $T^p$; summing
$-T\phi(u)^n\phi'(u)$ over $u\in k$ leaves only positive exponents
divisible by $p-1$.  This proves (8.16), and substitution yields (8.13).

Every admissible label in (8.11) satisfies $j+\nu\geq N_0-p$.  Together
with the displayed degrees of $B^0_{d,c}$, this proves the first equality
in (8.14) by the nilpotence threshold in
$k[x,y]/(x^{N_0},y^{N_0})$; the shears below preserve this filtration.

For the Cartier term first take $d=0$ and work in
\[
 \mathscr A=k[G,Z]/(G^{N_0}-1,Z^p-1),
 \qquad E=G^p,\quad T=E-1=x^p,\quad z=Z-1.
\]
The root set of
$Q(U)=\prod_{b\in k}(U-E^b)$ is stable under multiplication by $E$, so
$Q(EU)=Q(U)$.  If $q_i$ is an intermediate coefficient, then
$(E^i-1)q_i=0$; because $(E^i-1)/T$ is a unit for $1\leq i<p$, each $q_i$
is divisible by $T^{p-1}$.  The constant coefficient is $-1$, and
$Q(1)=0$.  Setting $U=Z$ and using $z^p=0$ therefore gives
\[
 \prod_{b\in k}(Z-G^{pb})\in x^{p(p-1)}z\mathscr A.
\]
Among $a=0,\ldots,N_0-p$ there remain $(p-1)^2$ indices not divisible by
$p$, and each corresponding factor $Z-G^a$ lies in $(x,z)$.  A monomial
surviving their product with the preceding ideal would have the form
\[
 x^{p(p-1)+\alpha}z^{1+\gamma},\qquad
 \alpha\leq p-1,\quad\gamma\leq p-2,\quad
 \alpha+\gamma\geq(p-1)^2.
\]
This is impossible for odd $p$, since $(p-1)^2>2p-3$.  Hence, with
$j_0=N_0-p+1$,
\[
 P_{j_0}(G,Z)=\prod_{a=0}^{N_0-p}(Z-G^a)=0
 \quad\text{in }\mathscr A,
 \tag{8.17}
\]
and therefore $P_j(G,Z)=0$ for every $j\geq j_0$.

Now $\rho_p(\Pi_0b_{j,\nu})$ is $x^\nu$ times a polynomial in $H^p$.
After multiplication by a nonzero $B^0_{0,c}$, every positive power of
$H^p-1=y^p$ exceeds the $y^{N_0-1}$ boundary; only evaluation at $H^p=1$
can survive.  Its coefficient is the image of $P_j(G,Z)$ modulo $Z^p-1$,
which vanishes for $j\geq j_0$ by (8.17).  If $j<j_0$, then
\[
 \beta_p(j)\leq\beta_p(N_0-p)=2p^2-3p+1,\qquad
 \nu\geq e-\beta_p(j)\geq(p-1)^3\geq N_0-p;
\]
the $x$-power in every nonzero auxiliary seed makes the product zero.
This proves the second equality of (8.14) for $d=0$.

For general $d$, let
$\Theta_db(X,Y)=b(X,X^dY)$ and let
$\Phi_d(r,s)=(r+ds,s)$.  The first map preserves $K_{\mathcal A}$ and
$\Pi_0\Theta_d=\Theta_d\Pi_0$, while
\[
 \rho_p(\Theta_db)=\Phi_d^{-1}(\rho_p(b)),
 \qquad B^0_{d,c}=\Phi_d(B^0_{0,c}).
\]
Conjugating the $d=0$ identity proves (8.14) for every residue.  Substitution
in the packet defect identity now gives (8.12).
\end{proof}

\begin{lemma}[Carrier vanishing]
\label[lemma]{lem:carrier-vanishing}
For every prime $p\geq5$ and every $d\in k$,
\[
 \mathcal H_d(K_{\mathcal A})\subset
 \mathcal C_{2,\kappa(d)}.
 \tag{8.18}
\]
In particular,
$\mathcal H_0(K_{\mathcal A})\subset\mathcal C_{2,\infty}$ and
$\mathcal H_A(K_{\mathcal A})\subset\mathcal C_2$.
\end{lemma}

\begin{proof}
For $p\geq5$, every label in (8.11) satisfies the sharper bound
$j+\nu\geq2p^2-p-1$.  Since the least total degree of a nonzero seed in
(8.13) is $p$, \cref{lem:seed-annihilation} gives
$\lambda_{d,c}(b_{j,\nu})=0$ for all $c$.  The mixed differences generate
$I_{\kappa(d)}$, so the Frobenius pairing in (8.2) identifies their joint
kernel with $I_{\kappa(d)}^\perp=\mathcal C_{2,\kappa(d)}$.  Linearity on
the Newton basis proves (8.18).
\end{proof}

\begin{lemma}[Second-domain cancellation]
\label[lemma]{lem:second-domain-cancellation}
In the recursive top-shell bases write
\[
 B_3^\transpose=
 \begin{pmatrix}B_2^\transpose&K\\0&Z\end{pmatrix},
 \qquad
 T=\begin{pmatrix}\mathcal A&0\\\mathcal X&pV\end{pmatrix}.
\]
For $f\in\mathcal J_2$, put $r_f=K^\transpose f$.  Then
$T^\transpose r_f$ is divisible by $p$, and, if
\[
 s_f=p^{-1}T^\transpose r_f\pmod p=(s_{f,1},s_{f,2})
\]
according to the two domain blocks of $T$, one has $s_{f,2}=0$.
\end{lemma}

\begin{proof}
For $(a,b)\in k^2$, let $F_{a,b}$ be the indicator of the fibre of
$G_2\to k^2$ above $(a,b)$ and put
\[
 Y_b=\sum_{a\in k}F_{a,b},\qquad
 \mathsf H_b=Y_b-F_{0,b},\qquad
 U=\sum_{b\ne0}Y_b.
\]
For a fine point $z$ above $(a,b)$, the sum of one line through $z$ in
each projective direction is $\mathbf1-F_{a,b}$ modulo $p$.  Thus every
$F_{a,b}$, and hence $Y_b$, $\mathsf H_b$ and $U$, belongs to the
depth-two line code $\mathcal C_2$.

The coefficient sum for the first domain block is
\[
 \overline{KT_1}(X^iY^j)
 =\mathbf1_{\{\bar j=0\}}U
  +\mathbf1_{\{\bar i=0\}}\mathsf H_{\bar j}.
\]
Pairing with $f\in\mathcal J_2=\mathcal C_2^\perp$ proves the required
divisibility on that block.  After removing the explicit factor $p$ from
the second diagonal block, its coordinate indexed by
$0\leq r<N_0$, $0\leq i<e$ has point kernel
\[
 D_{r,i}=\mathbf1_{\{r\equiv0\pmod p\}}\mathsf H_{\bar i}.
\]
Indeed the low-digit equation in its defining coefficient packet is
\[
 ps(x-r)\equiv i-y\pmod{p^2}.
\]
If $x-r$ is a nonunit the packet cancels; if it is a unit there are $p$
lifts, whose sum is zero unless $r\equiv0\pmod p$, and in the remaining
case the packet is exactly $\mathsf H_{\bar i}$.  Since this function is
in $\mathcal C_2$, every such coordinate pairs to zero with $f$.  This
proves $s_{f,2}=0$ and the lemma.
\end{proof}

\begin{proposition}[Exact divided-transfer sequence]
\label[proposition]{prop:divided-transfer-sequence}
For every prime $p$, let $\pi_*:\A_3\to\A_2$ sum coefficients on point fibres and
let
\[
 \mathscr H_3^{(2)}=
 \{\widetilde y\in\A_3/p^2\A_3:B_3\widetilde y=0\},
 \qquad
 \mathcal J_2:=\mathcal C_2^\perp
 =\bigcap_{\delta\in\mathbb P^1(k)}I_\delta.
\]
Put
\[
 S_3=(\LL_3:p),\qquad
 \overline S_3=(S_3+p\A_3)/p\A_3,
 \qquad \mathcal Y_3=\overline S_3^\perp.
\]
Then
\[
 \widetilde{\mathcal C}_3:\mathscr H_3^{(2)}\longrightarrow\mathcal J_2,
 \qquad \widetilde y\longmapsto p^{-1}\pi_*\widetilde y\pmod p
\]
factors through reduction onto $\mathcal Y_3$, and the induced map
$\mathcal C_3:\mathcal Y_3\to\mathcal J_2$ is the transpose of $\eta_2$
under the natural finite-field pairings.  If
$\mathcal X_A=\mathcal X|_{K_{\mathcal A}}$ and
$\mathcal Q_p:=\coker(\mathcal X_A^\vee)$, there is a canonical map
\[
 \Theta:\mathcal J_2\longrightarrow\mathcal Q_p,
 \qquad f\longmapsto[\Lambda_f],\qquad
 \Lambda_f(b)=\langle s_{f,1},b\rangle
 \quad(b\in K_{\mathcal A}),
\]
for which
\[
 \mathcal Y_3\xrightarrow{\ \mathcal C_3\ }
 \mathcal J_2\xrightarrow{\ \Theta\ }\mathcal Q_p
 \quad\text{is exact at }\mathcal J_2.
\tag{8.19}
\]
Furthermore, dualizing the kernel sequence of $\mathcal X_A$ gives
\[
 \mathcal Q_p\cong K_{\mathcal A\mathcal X}^\vee,
 \qquad [\Lambda_f]\longmapsto
 \Lambda_f|_{K_{\mathcal A\mathcal X}}.
\tag{8.20}
\]
\end{proposition}

\begin{proof}
The source sequence and its finite-field dual are
\[
 0\longrightarrow\ker\eta_2\longrightarrow\QQ_2/p\QQ_2
 \xrightarrow{\ \eta_2\ }p\QQ_3/p^2\QQ_3,
\]
\[
 (p\QQ_3/p^2\QQ_3)^\vee\xrightarrow{\ \eta_2^\vee\ }
 \mathcal J_2\longrightarrow(\ker\eta_2)^\vee\longrightarrow0,
\tag{8.19a}
\]
where $\mathcal J_2\cong(\QQ_2/p\QQ_2)^\vee$.  To realize the first
dual space, reduction modulo $p$ maps
$\mathscr H_3^{(2)}$ onto $\mathcal Y_3\cong
(p\QQ_3/p^2\QQ_3)^\vee$.  Indeed, reduction of a mod-$p^2$ harmonic
annihilates both $\LL_3/p\LL_3$ and every divided relation.  Conversely,
if $y\in\mathcal Y_3$, then $B_3y=pu$; orthogonality to the divided
relations gives $u\in\im\overline B_3$, so a choice of
$z$ with $\overline B_3z=-u$ makes $y+pz$ harmonic modulo $p^2$.

By \cref{cor:first-transfer-congruence}, for every depth-two point basis
vector $a$ one may write $\tau_3(a)=B_3^\transpose c+px$.  Adjunction of
fibre summation and transfer shows that $\pi_*\widetilde y$ is divisible by
$p$ modulo $p^2$ for $\widetilde y\in\mathscr H_3^{(2)}$.  If two such
lifts differ by $pz$ with $B_3z=0$ modulo $p$, then
$\langle\pi_*z,a\rangle=\langle z,\tau_3(a)\rangle=0$ modulo $p$ by the
same congruence; hence the raw map depends only on the reduction modulo
$p$.  Pairing $\widetilde{\mathcal C}_3$ with a depth-two line and using
the depth-transition identity puts its image in
$\mathcal J_2=\mathcal C_2^\perp$.  Finally,
\[
\left\langle\widetilde{\mathcal C}_3(\widetilde y),a\right\rangle
 =\frac{\langle\widetilde y,\tau_3(a)\rangle}{p}
 \equiv\langle y,x\rangle\pmod p,
\]
which is precisely the dual Bockstein defining $\eta_2^\vee$.
Consequently $\widetilde{\mathcal C}_3$ factors through reduction onto
$\mathcal Y_3$, and the induced map $\mathcal C_3$ is the transpose of
$\eta_2$ in (8.19a).

For exactness, keep the recursive bases of
\cref{lem:second-domain-cancellation}.  Put
$R_{\mathrm{top}}=(N+N_0)e$.  The complementary top-shell blocks satisfy
\[
 Z^\transpose T^\transpose=T^\transpose Z^\transpose
 =p^3I_{R_{\mathrm{top}}}
\]
(the $n=3$ specialization of the integral dual identity in
\cite[Theorem~8.3, equation~(26)]{BAB26D2}).  A target
$f\in\mathcal J_2$ lies in
$\im\mathcal C_3$ exactly when the prescribed coarse point component
$pf$ extends to a mod-$p^2$ harmonic.  Transposing the graph form in
\cref{lem:second-domain-cancellation} gives
$B_3=\left(\begin{smallmatrix}B_2&0\\K^\transpose&Z^\transpose\end{smallmatrix}\right)$.
The first block is soluble because $f\in\mathcal J_2$; the top-shell block
is precisely the requirement that
\[
 Z^\transpose v\equiv-pK^\transpose f\pmod{p^2}
\]
be solvable.  Put $r_f=K^\transpose f$; replacing $v$ by $-v$ removes
the harmless sign.  Multiplying by $T^\transpose$,
one direction gives
\[
 T^\transpose r_f\in
 pT^\transpose\Z_p^{R_{\mathrm{top}}}
 +p^2\Z_p^{R_{\mathrm{top}}}.
\tag{8.19b}
\]
Conversely, if
$T^\transpose r_f=pT^\transpose w+p^2v$, multiplication by
$Z^\transpose$ gives
$p^3(r_f-pw)=p^2Z^\transpose v$; cancellation in the torsion-free
$p$-adic lattice yields
$pr_f=Z^\transpose v+p^2w$.  Hence (8.19b) is equivalent to the original
congruence, in both directions.

The two-chart reduction is
\[
 \overline T=\begin{pmatrix}\mathcal A&0\\\mathcal X&0\end{pmatrix}.
\]
By \cref{lem:second-domain-cancellation}, the divided adjoint
$s_f=(s_{f,1},0)$.  Therefore $s_f\in\im\overline T^\transpose$ exactly
when $s_{f,1}$ lies in the combined row space of $\mathcal A$ and
$\mathcal X$.  Restriction to $K_{\mathcal A}$ kills the first-chart row
space, so the remaining condition is exactly
$\Lambda_f\in\im\mathcal X_A^\vee$.  For
$b\in K_{\mathcal A\mathcal X}$, expanding $s_f$ in the fixed cyclotomic
bases gives the coefficient formula (8.5), hence
\[
 \Lambda_f(b)=\langle f,\delta_1(b)\rangle.
\]
Thus $\Lambda_f$ is defined on all of $K_{\mathcal A}$ before restriction,
whereas $\delta_1$ is used only on its typed domain
$K_{\mathcal A\mathcal X}$; no unnamed section is being used.  This proves
exactness in (8.19).

Finally, dualizing
$0\to K_{\mathcal A\mathcal X}\to K_{\mathcal A}
\xrightarrow{\mathcal X_A}\im\mathcal X_A\to0$
gives (8.20).  The argument uses only the integral incidence lattices, the
first transfer congruence, the top-shell dual identity and finite-dimensional
duality; no Smith multiplicity or conclusion of \cref{thm:main-profile}
enters.
\end{proof}

\begin{lemma}[Divided-trace duality]
\label[lemma]{lem:divided-trace}
For $p\geq5$, the following are equivalent:
\[
 d_p=0;\qquad \delta_1(K_{\mathcal A\mathcal X})=0;\qquad
 \mathcal H_X(K_{\mathcal A\mathcal X})\subset\mathcal C_2.
\tag{8.21}
\]
\end{lemma}

\begin{proof}
Since $\mathcal C_3=\eta_2^\vee$, one has $d_p=0$ exactly when
$\mathcal C_3$ is onto.  Exactness in
\cref{prop:divided-transfer-sequence} makes this equivalent to
$[\Lambda_f]=0$ for all $f\in\mathcal J_2$.  Under (8.20), this says
\[
 \langle f,\delta_1(b)\rangle=0
 \quad(f\in\mathcal J_2, b\in K_{\mathcal A\mathcal X}).
\]
The perfect pairing
$\mathcal J_2\times(\mathscr R_2/\mathcal C_2)\to k$ makes this equivalent
to $\delta_1(K_{\mathcal A\mathcal X})=0$.
\Cref{lem:carrier-vanishing} gives
$\mathcal H_A(K_{\mathcal A\mathcal X})\subset\mathcal C_2$, so (8.6)
gives the final equivalence.
\end{proof}

\begin{theorem}[Cross-chart divided-carry theorem]
\label[theorem]{thm:o2}
For every prime $p$, $d_p=0$; equivalently, (O2) holds.
\end{theorem}

\begin{proof}
Let $p\geq5$ and $b\in K_{\mathcal A\mathcal X}$.  By
\cref{lem:carry-transpose}, $Sb\in K_{\mathcal A}$, and
\cref{lem:carrier-vanishing} puts $\mathcal H_0(Sb)$ in
$\mathcal C_{2,\infty}$.  Equations (8.8)--(8.9) therefore give
$\mathcal H_X(b)\in\mathcal C_{2,0}\subset\mathcal C_2$.
\Cref{lem:divided-trace} yields $d_p=0$.

The primes $3$ and $2$ are not inferred from this degree bound.  The exact
$p=3$ double-kernel certificate checks all $360$ zero-cross Newton labels
(the $18$ candidate columns have rank $18$), hence gives
$\eps_1(3,3)=18$.  At $p=2$ the semisimple sector and division by two are
unavailable.  A separate full-matrix certificate constructs $\mathcal A$ and
$\mathcal AS$ directly and finds ranks $8$ and $9$, so (4.2) gives
$\eps_1(2,3)=1$ without using a Smith multiplicity.  For both primes, two
independent Smith lanes give the complete profile; combining those rows with
the modular-rank values $a_0$ and these independently obtained obstruction
values in the defect identity gives $d_2=d_3=0$.  The distributed $p=3$
cross-chart certificate additionally checks the divided-carry membership
directly; no separate direct O2-membership claim is made for $p=2$.  The
exact inputs, hashes and non-circular deductions are recorded in
\cref{app:exceptional-branches}.
\end{proof}

The proof uses the exchange only in the direction (8.7); it does not claim
that coordinate swap is an involution on the fixed two-chart splitting.
Likewise, (8.8) is an identity after compatible candidate relations are
combined, not a columnwise divided identity.

\section{Exact assembly and the all-prime Smith profile}
\label{sec:assembly}

Put
\[
 N_3=p^6,\qquad
 \sigma_p=\vp|\QQ_3|=\frac{3p^6-p^5-p^3-p}{2},
\tag{9.1}
\]
and retain $\beta,\gamma$ from (2.5).  By
\cref{lem:defect,thm:o2},
\[
 T_{\geq3}:=m_3+m_4+m_5=a_0+\eps_1(p,3).
\tag{9.2}
\]
The symbol $T_{\geq3}$ is used only for this tail sum; $H_p$ remains the
boundary-column count from (6.5).

Rank, total size and weighted order give
\begin{align}
 m_0&=M_p,\nonumber\\
 m_1&=2N_3-2M_p+T_{\geq3}+\beta+2\gamma-\sigma_p,\nonumber\\
 m_2&=\sigma_p-N_3+M_p-2T_{\geq3}-\beta-2\gamma,
 \tag{9.3}\label{eq:assembly}\\
 m_3&=T_{\geq3}-\beta-\gamma,\nonumber\\
 m_4&=\beta,\qquad m_5=\gamma.\nonumber
\end{align}

\begin{theorem}[All-prime depth-three Smith profile]
\label[theorem]{thm:main-profile}
For $p=2$ and $p=3$, respectively,
\[
 (m_0,m_1,m_2,m_3,m_4,m_5)=\FCIGDThreeSmithProfileTwo,
\tag{9.4}
\]
\[
 (m_0,m_1,m_2,m_3,m_4,m_5)=\FCIGDThreeSmithProfileThree.
\tag{9.5}
\]
For every prime $p>3$ with $p\equiv1\pmod6$,
\[
\FCIGDThreeSmithProfileModOne
\tag{9.6}
\]
For every prime $p>3$ with $p\equiv5\pmod6$,
\[
\FCIGDThreeSmithProfileModFive
\tag{9.7}
\]
Equivalently,
\[
 (\QQ_3)_{(p)}\cong
 \bigoplus_{h=1}^{5}(\Z/p^h\Z)^{m_h}.
\tag{9.8}
\]
\end{theorem}

\begin{proof}
The transfer theorem gives $m_4,m_5$ and excludes higher exponents.
\Cref{thm:modrank} gives $m_0$, \cref{prop:a0} gives $a_0$,
\cref{thm:epsilon} gives $\eps_1$ for odd primes, and the exact exceptional
certificate gives $\eps_1(2,3)=1$.  By \cref{thm:o2}, the defect term in
\cref{lem:defect} is zero.  All hypotheses of \eqref{eq:assembly} therefore
hold for every prime.  Direct substitution gives (9.4)--(9.7).  No value at
$p=5,7$ or $11$ is used to select either branch.
\end{proof}

The formulas satisfy identically
\[
 \sum_{h=0}^{5}m_h=p^6,
 \qquad
 \sum_{h=1}^{5}hm_h=\frac{3p^6-p^5-p^3-p}{2}.
\tag{9.9}
\]
After substituting $p=6h+1$ in (9.6), every multiplicity expands as a
polynomial in $\Z_{\geq0}[h]$; for example
\[
 m_2=3h(4h+1)(6h+1)
 (216h^3+102h^2+25h+2),
\]
and the same direct expansion of (9.7) after $p=6h+5$ lies in
$\Z_{\geq0}[h]$ for all six multiplicities.  The public companion performs
both substitutions with exact rational arithmetic and rejects a nonintegral
or negative coefficient.  Thus integrality and positivity do not rely on a
bounded prime sweep.

\section{Finite controls, source boundary, and limitations}
\label{sec:controls}

The first profile rows are
\[
\begin{array}{c|rrrrrr}
p&m_0&m_1&m_2&m_3&m_4&m_5\\ \hline
2&30&6&19&6&2&1\\
3&240&168&195&108&15&3\\
5&4050&4140&4795&2490&140&10\\
7&26740&34608&36225&19488&567&21\\
11&363000&545622&561979&297330&3575&55
\end{array}
\tag{10.1}
\]
The $p=2,3$ rows are reconstructed independently by the public exceptional
packages and also agree with the assembled formulas.  The $p=5,7,11$ rows
are direct exact evaluations of (9.6)--(9.7) and of the tracked structured
formula object.  All five rows are public regression controls; none is an
input used to select a residue-class branch.

The active common-image thresholds at the three hostile odd-prime controls
are
\[
\begin{array}{c|l}
p&\{\tau_p(K)<e\}\\ \hline
3&9,12,15\\
5&45,50,55,60,64,88,92,96\\
7&114,120,126,133,140,147,154,203,210,216,222,228,234,240,288.
\end{array}
\tag{10.2}
\]
At $p=5$, the threshold $64$ is a required pattern break: any rule that
predicts only an arithmetic progression fails there.

The bounded literature record is not novelty clearance.  Molchanov treats
the affine residue-ring Radon transform, its normal operator, spectrum and
inversion \cite{MOL12}; Ilmavirta gives the finite-group transform context
\cite{ILM}.  Adjacent projective Hjelmslev rank and lifting results are due
to Landjev--Vandendriessche, Kohnert and Hanssens--Van Maldeghem
\cite{LV14,KOH09,HVM89}.  \L aba--Trainor and Dvir study the modular
hyperplane-function object and rank bounds \cite{LT25,DVI26}, not the
integral affine Smith profile here.  Two additional projective-rank sources
were available only at metadata or abstract level \cite{LAN17,BLR23}; no
theorem is transferred from them.

The present proof is depth-specific.  It does not prove a standard-basis or
Rees strictness theorem in general, an all-depth recurrence for closed Smith
profiles, or a complexity improvement for the full incidence matrix.  The
Newton-minors argument succeeds because it computes the actual truncated
image directly; it does not justify the false identity
$\operatorname{in}(U\cap V)=\operatorname{in}(U)\cap
\operatorname{in}(V)$ in arbitrary filtered modules.

\section{Public reproducibility
\texorpdfstring{\protect\\}{ }and companion interface}
\label{sec:repro}

The public repository contains the cited-only manuscript source, the
structured formula object from which the displayed branches are generated,
an exact-arithmetic companion, bounded certificates for the exceptional
primes, hostile tests, and fail-closed source and release manifests.  The
one-command verifier checks:
\begin{enumerate}
\item the literal formulas (7.2), (7.7), and (9.4)--(9.7);
\item both polynomial branches of (9.9), integrality substitutions and the
  finite rows (10.1);
\item the typed distinction among vector $G_K$, boundary count $H_p$,
  generator invariant $\ngen_{\overline\OO}(W)$ and tail sum
  $T_{\geq3}$;
\item the symbolic rank and weighted-order identities for both residue
  classes modulo $6$, together with independent numerical controls;
\item the exact $p=2$ two-chart rank gap $9-8=1$ and its restriction to
  the $24$-dimensional first-chart kernel;
\item source-manifest, PDF, current-checkout and release-surface integrity.
\end{enumerate}

The layout follows the Depth-2 companion \cite{BAB26D2}:
\texttt{paper/} contains the source, generated formula fragment and
title-named PDF; \texttt{companion/} contains the structured formulas and
exact arithmetic; \texttt{tests/} contains mutation-sensitive controls;
\texttt{scripts/} supplies the verifier; and \texttt{docs/} states source,
claim and reproducibility boundaries.  Reproduction commands and the
hash-locked dependency environment are documented in
\texttt{REPRODUCE.md}.  Finite replay is a supporting check and is not
used as a proof of the all-prime theorem.

\appendix

\section{Typed symbol contract}
\label{app:symbols}

The following roles are normative:
\[
\begin{array}{c|l}
N=p^3,\ N_0=p^2&\text{depth-three and coarse residue moduli}\\
E_{\mathcal A},E_{\mathcal X}&\text{first-chart and cross-chart targets}\\
\mathfrak q&\text{formal Gaussian-polynomial indeterminate}\\
G_K&\text{boundary vector in }\overline\OO^{p^2}\\
H_p&\text{number of admitted boundary columns}\\
\tau_p(K)&\text{DVR Smith threshold}\\
\ngen_{\overline\OO}(W)&\text{minimal number of module generators}\\
T_{\geq3}&m_3+m_4+m_5\\
\mathcal H_d,\mathcal H_X&\text{first-chart residue and cross divided carries}\\
\mathcal C_{2,\delta}&\text{depth-two tangent-bundle line code}.
\end{array}
\]
In particular, no scalar $G_p$ is defined, and the auxiliary letter $H$ is
not used for the Smith tail.

\section{Residue-block arithmetic for the closed sum}
\label{app:arithmetic-blocks}

This appendix supplies the complete block decomposition used in
\cref{lem:tail-endpoints}.  If $M=ap+b$ with $0\leq b<p$, then
\[
 \sum_{K=0}^{M-1}\left\lfloor\frac Kp\right\rfloor
 =p\frac{a(a-1)}2+ab.
\tag{A.1}
\]
For $p=6h+1$ one has
\[
 M=2hp+h=h(2p-1)+2h,\qquad c=4h.
\]
Thus (A.1), with $(a,b)=(2h,h)$, gives the first floor sum.  For the
shifted denominator $d=2p-1=12h+1$, the interval
\[
 c\leq K+c\leq c+M-1=hd+6h-1
\]
has an initial block of length $d-c=8h+1$, $h-1$ complete intermediate
blocks, and a terminal block of length $6h$.  Hence
\[
 \sum_{K=0}^{M-1}\left\lfloor\frac{K+c}{d}\right\rfloor
 =d\sum_{j=1}^{h-1}j+6h^2
 =d\frac{h(h-1)}2+6h^2.
\tag{A.2}
\]
The ramp in a complete $d$-block is
\[
 \sum_{r=p}^{d-1}(r-(p-1))=\frac{p(p-1)}2
 =3h(6h+1).
\]
It occurs in the initial block and the $h-1$ complete blocks, but not in
the terminal block, giving $\sum\delta_K=3h^2p$.

For $p=6h+5$,
\[
 M=(2h+1)p+4h+3,\qquad d=12h+9,\qquad c=7h+5.
\]
Equation (A.1) with $(a,b)=(2h+1,4h+3)$ gives the first floor sum.  Now
\[
 c+M-1=(h+1)d+6h+3.
\]
There is an initial block of length $d-c=5h+4$, followed by $h$ complete
blocks and a terminal block of length $6h+4=p-1$.  Consequently
\[
 \sum_{K=0}^{M-1}\left\lfloor\frac{K+c}{d}\right\rfloor
 =d\frac{h(h+1)}2+(h+1)(p-1).
\tag{A.3}
\]
The initial ramp omits the values $1,\ldots,h$, each complete block
contributes $p(p-1)/2$, and the terminal ramp is zero.  Therefore
\[
 \sum_{K=0}^{M-1}\delta_K
 =(h+1)\frac{p(p-1)}2-\frac{h(h+1)}2,
\tag{A.4}
\]
which is the last entry of (7.9).  These decompositions also give
$L_{M-1}=p-1$ and $L_M=0$ directly.

\section{Exceptional branches and executable certificates}
\label{app:exceptional-branches}

The primes $2$ and $3$ are separate proofs, not data used to fit the
odd-prime formulas.  The executable package is
\begin{center}
\path{companion/exceptional_profiles/}.
\end{center}
It constructs the canonical point-by-line transpose $B_3^\transpose$ and
computes its local Smith valuations by two independent exact lanes:

\begin{enumerate}
\item minimum-valuation pivot elimination over $\Z/p^7\Z$;
\item Bockstein layer peeling.
\end{enumerate}

Both lanes return
\[
 \begin{array}{c|c|c|c}
p&\operatorname{shape}(B_3^\transpose)&(m_0,\ldots,m_5)&
 \sum_{h=1}^{5}h\,m_h\\ \hline
 2&64\times96&\FCIGDThreeSmithProfileTwo&75\\
 3&729\times972&\FCIGDThreeSmithProfileThree&957.
 \end{array}
\tag{B.1}
\]
The matrix SHA-256 values are respectively
\[
 \begin{split}
 &\texttt{CAD9CC122B4811426CE958D1EA9487C536BF51FEACD3550FA04A94A012EB5E82},\\
 &\texttt{3D05CC493F837A76693934970005F025F3CC37C70DCECDD659295A37BC23CC68}.
 \end{split}
\]
The canonical result and verifier source are authenticated by the public
source manifest; normal and optimized executions reproduce the same result
bytes.

The same directory contains a separate first-transfer certificate.  At
$p=2$ and $p=3$ it solves exactly for one coarse point-fibre indicator in
the modular depth-three line code, checks the solution residual, and then
independently pairs the indicator with a basis of the full dual nullspace.
The matrix ranks are $30$ and $240$, the nullities are $34$ and $489$, and
the displayed primal certificates use $14$ and $147$ lines, respectively.
Deleted-line and flipped-point mutations are rejected.  Affine translation
therefore proves the first-transfer congruence for every point fibre at the
two exceptional primes.

The independent package
\begin{center}
\path{companion/epsilon1_p2/}
\end{center}
certifies the obstruction in (4.2) at $p=2$ without reading (B.1).  In the
power basis of $\F_2[Z]/(Z^4+1)$, its $32$ columns are indexed by
$X^iY^j$, $0\leq i<4$, $0\leq j<8$, and its rows by
$(t,b)$, $0\leq t<8$, $0\leq b<4$.  It constructs
\[
 \mathcal A_{(t,b),(i,j)}=[Z^b]Z^{i+tj},\qquad
 (\mathcal AS)_{(t,b),(i,j)}=[Z^b]Z^{ti+j}.
\tag{B.2}
\]
Two independent exact rank algorithms give
\[
 \rank_{\F_2}\mathcal A=8,\qquad
 \rank_{\F_2}\begin{bmatrix}\mathcal A\\ \mathcal AS\end{bmatrix}=9,
 \qquad \dim_{\F_2}K_{\mathcal A}=24.
\tag{B.3}
\]
Direct restriction gives
$\rank((\mathcal AS)|_{K_{\mathcal A}})=1$; unit-slope rows contribute
rank zero on the kernel, while nonunit-slope rows already contribute rank
one.  The certificate also records an invertible $9\times9$ witness minor
and checks every cyclotomic remainder by both a closed formula and polynomial
long division.  Thus (4.2) yields $\eps_1(2,3)=1$ independently of the Smith
row.

There is no circular use of \cref{thm:main-profile}: (B.1) is obtained
directly from the incidence matrix.  In the defect identity (8.1), the tails
from (B.1) are
\[
 6+2+1=9=8+\eps_1(2,3),\qquad
 108+15+3=126=108+\eps_1(3,3),
\]
where the independently certified obstruction values are $1$ from
(B.2)--(B.3) and $18$ from the $p=3$ cross-chart rank.
Thus the direct rows imply $d_2=d_3=0$.  For $p=3$ the cross-chart package
in \path{companion/cross_chart_p3/} reconstructs all $378$ Newton labels,
classifies $360$ individually zero cross-chart labels, proves that the
remaining $162\times18$ matrix has rank $18$, and checks every divided
carry against both the correct line-row space and its annihilator
differences.  This is an additional local check, not the source of (B.1).

From the repository root, the replay commands are to be run sequentially:
\begin{verbatim}
python companion/exceptional_profiles/verify_exceptional_profiles.py
python -O companion/exceptional_profiles/verify_exceptional_profiles.py
python companion/exceptional_profiles/verify_o1_exceptional.py
python -O companion/exceptional_profiles/verify_o1_exceptional.py
python companion/cross_chart_p3/verify_cross_chart_p3.py
python -O companion/cross_chart_p3/verify_cross_chart_p3.py
python companion/epsilon1_p2/verify_epsilon1_p2.py
python -O companion/epsilon1_p2/verify_epsilon1_p2.py
\end{verbatim}
Each verifier is read-only and authenticates its result and execution
receipt.  No exceptional computation is extrapolated to another prime.

\section*{AI-use disclosure and author responsibility}

Under the author's direction, OpenAI Codex (GPT-5.6 Sol) generated the central
mathematical development---including the all-prime depth-three Smith profile,
Newton-minor analysis and divided-carry closure---and assisted with software,
certificates and manuscript.  The author checked the work, set its public
boundaries and takes responsibility for the final content.  Same-workflow
machine reviews are not human peer review or independent expert verification.
Detailed provenance is recorded in the repository's \path{AI_USE.md}.

\printbibliography

\end{document}
